% concatenated from Sanjay_Mallick.tex
%Please name your file by your name.

\documentclass[a4paper,draft]{amsproc}
\usepackage{amssymb}
\usepackage[hyphens]{url} \urlstyle{same}
%\usepackage[dvips]{graphicx} 
%\usepackage{refcheck}
%\usepackage{cmdtrack} %% Checks whether all author defined macros are used or not0
\theoremstyle{plain}
\newtheorem*{theoA}{Theorem A}
\newtheorem*{theoB}{Theorem B}
\newtheorem*{theoC}{Theorem C}
\newtheorem*{theoD}{Theorem D}
\newtheorem*{theoE}{Theorem E}
\newtheorem*{theoF}{Theorem F}
\newtheorem{theo}{Theorem}[section]
%\newtheorem{prop}{Proposition}[section]
\newtheorem{lem}{Lemma}[section]
\newtheorem{cor}{Corollary}[section]
\newtheorem{exm}{Example}[section]
\newtheorem{defi}{Definition}[section]
\newtheorem{ques}{Question}[section]
\theoremstyle{definition}
%\newtheorem{dfn}{Definition}[section]
\theoremstyle{remark}
\newtheorem{rem}{Remark}[section]

\newcommand{\ol}{\overline}
\newcommand{\be}{\begin{equation}}
	\newcommand{\ee}{\end{equation}}
\newcommand{\beas}{\begin{eqnarray*}}
	\newcommand{\eeas}{\end{eqnarray*}}
\newcommand{\bea}{\begin{eqnarray}}
	\newcommand{\eea}{\end{eqnarray}}
\renewcommand{\d}{\displaystyle}
\newtheorem*{ack}{Acknowledgement}
\numberwithin{equation}{section}

%% Please, do not change the following four lines:
\renewcommand{\le}{\leqslant}\renewcommand{\ge}{\geqslant}
\renewcommand{\leq}{\leqslant}\renewcommand{\geq}{\geqslant}
\renewcommand{\setminus}{\smallsetminus}
\setlength{\textwidth}{28cc} \setlength{\textheight}{42cc}

\title[A GENERAL CHARACTERIZATION ON THE UNIQUENESS PROBLEM...]{A GENERAL CHARACTERIZATION ON THE UNIQUENESS PROBLEM OF L-FUNCTIONS AND GENERAL MEROMORPHIC FUNCTIONS}

\subjclass[2020]{11M36; 30D35}

\keywords{L-function; meromorphic function; shared set; uniqueness}

\author[Mallick]{\bfseries Sanjay Mallick}
\address{
	Department of Mathematics \\ % \hfill (Received 00 00 202?)\\
	Cooch Behar Panchanan Barma University   \\ %\hfill (Revised  00 00 202?)\\
	Cooch Behar\\
	India}
\email{sanjay.mallick1986@gmail.com, smallick.ku@gmail.com}

%% OTHER AUTHOR(S):
\author[Saha]{\bfseries Ripan Saha}
\address{ 	Department of Mathematics \\ % \hfill (Received 00 00 202?)\\
	Cooch Behar Panchanan Barma University   \\ %\hfill (Revised  00 00 202?)\\
	Cooch Behar\\
	India }
\email{ripan.tfg.saha@gmail.com}

\thanks{Supported by SERB, DST, INDIA and NBHM, DAE, INDIA} 
\thanks{Communicated by ...}

\begin{document}
	
	{\begin{flushleft}\baselineskip9pt\scriptsize
			%PUBLICATIONS DE L'INSTITUT MATH\'EMATIQUE\\
			%Nouvelle s\'erie, tome 1??(1??) (202?), od--do \hfill DOI: \\
			MANUSCRIPT
	\end{flushleft}}
	\vspace{18mm} \setcounter{page}{1} \thispagestyle{empty}
	
	
	\begin{abstract}
		 In the paper, concerning a question of Yi \cite{yuan2018}, we study general criterion for the uniqueness of an L-function and a general meromorphic function. Our results improve and extend all the existing results in this direction \cite{yuan2018,ss,sh,bk} to the most general setting. Moreover, we have exhibited a handsome number of examples to justify our claims as well as to confirm the wide-ranging applications of our results.  
	\end{abstract}
	
	\maketitle
	
	\section{Introduction and Related Results} 
	
	 The value distribution  of L-functions in the Selberg class via Nevanlinna theory is a recent trend to study different properties of L-functions as well as uniquely determining the same. This kind of L-functions includes the Riemann zeta function $\zeta(s) = \sum\limits_{n=1}^{\infty} n^{-s}$ and essentially those Dirichelet series where one might expect a Riemann hypothesis. Such an L-function is
	defined \cite{selberg1992,steuding2007value} to be a Dirichelet series
	\be\nonumber
	\mathcal{L}(s) = \sum\limits_{n=1}^{\infty} \frac{a(n)}{ n^{-s}}
	\ee
	with $a(1)=1$ satisfying the following axioms:
	\begin{itemize}
		\item(i) \emph{Ramanujan hypothesis :} $a(n)\ll n^{\varepsilon}$ for every $\varepsilon> 0;$
		\item(ii) \emph{Analytic continuation :} There is a non-negative integer $m$ such that\\ $(s-1)^m\mathcal{L}(s)$ is an entire function of
		finite order;
		\item (iii) \emph{Functional equation:} $\mathcal{L}$ satisfies a functional equation of type \begin{equation*}
			\Lambda_\mathcal{L}(s) = \omega\ol{ \Lambda_\mathcal{L}(1-\ol{s})},
		\end{equation*}
		where \begin{equation*}
			\Lambda_\mathcal{L}(s) = \mathcal{L}(s)Q^s\prod_{j=1}^{k}\Gamma(\lambda_js+\nu_j),
		\end{equation*}
		with positive real numbers $Q, \lambda_j$ and complex numbers $ \nu_j, \omega$ with Re$\nu_j \geq 0$ and $|\omega| = 1;$
		\item (iv) \emph{Euler product hypothesis :} $\log \mathcal{L}(s) =  \sum\limits_{n=1}^{\infty} \frac{b(n)}{ n^{s}}$,  where $b(n) = 0$ unless $n$ is a positive power of a prime
		and $b(n) \ll n^{\theta}$ for some $\theta < \frac{1}{2}$.
	\end{itemize}
	%All of the L-functions in the paper are assumed to be the Dirichelet series from the extended Selberg class only satisfying the axioms $(i)-(iii)$\cite{kaczorowski1999structure,steuding2007value}. 
	\par Throughout the paper by an L-function we always mean an L-function of the above kind and by any meromorphic function we always mean a meromorphic function defined in $\mathbb{C}$. We denote $\mathbb{\ol C}=\mathbb{C}\cup\{\infty\}$. By $\mathbb{N}$ we  mean the set of all natural numbers.  Though for standard definitions used in this paper we refer our readers to follow \cite{wkh}, yet for the sake of our convenience we denote the order of $f$ by $\rho(f)$, where 
	\be\nonumber\rho(f)=\limsup_ {r\rightarrow\infty} \frac{\log(T(r,f))}{\log r}.\ee
	By $S(r, f)$ we mean any quantity satisfying $S(r, f) = O(\log(rT(r, f)))$ for all $r$ possibly outside a set of finite linear measure. If $f$ is a function of finite order, then $S(r, f) = O(\log r)$
	for all $r$. 
	\vspace{0.1in}\par However, in recent times mainly during the last ten years,  a lot of investigations \cite{kac99,li2010,yuan2018,hu2016,selberg1992,steuding2007value}  have been pursued by various authors in the direction of value distribution of L-functions. In due course of time, this research has been moved towards uniquely determining the L-functions with respect to the meromorphic functions having finitely many poles so that we may know about various  properties of L-functions via the established literature of this special kind of meromorphic functions. Below we recall some of these results and their gradual developments. Before that, we need to address the following basic definitions to make out the results clearly. 
	% \\\\The importance of L-functions in number theory is needless to say and an L-function can be analytically continued to a meromorphic function in $\mathbb{C}$. Hence like the value distribution of meromorphic functions, the value distribution of L-functions is a natural consequence. In this respect, during the last few years an extensive study for the distribution of zeros of L-functions have been done by various researchers \cite{kac99,li2010,yuan2018,hu2016,selberg1992,steuding2007value}.    In due course of time, the study have been confined to the direction of uniquely determining an L-function via the shared values or sets. Hence let us recall these basic definitions of value and set sharing.
	\begin{defi}\cite{3a}
		For a  non-constant meromorphic function $f$ and  $a\in{\mathbb{C}}$, let  $E_{f}(a)=\{(z,p)\in\mathbb{C}\times\mathbb{N}: f(z)=a\; with\; multiplicity\; p\}$ \\$ \left(\ol  E_{f}(a)=\{(z,1)\in\mathbb{C}\times\mathbb{N}: f(z)=a\}\right)$, then we say  $f$, $g$ share the value $a$ CM(IM) if $E_{f}(a)=E_{g}(a)$$\left( \ol E_{f}(a)=\ol E_{g}(a)\right).$ For $a=\infty$, we define $E_f(\infty) := E_{1/f}(0)$ $\left( \ol E_f(\infty) :=\ol  E_{1/f}(0)\right)$.
	\end{defi}
	\begin{defi}\cite{3a}
		For a  non-constant meromorphic function $f$ and  $S\subset\overline{\mathbb{C}}$, let $E_{f}(S)=\bigcup_{a\in S}\{(z,p)\in\mathbb{C}\times\mathbb{N}: f(z)=a\; with\; multiplicity\; p\}$\\ $\left(\ol  E_{f}(S)=\bigcup_{a\in S}\{(z,1)\in\mathbb{C}\times\mathbb{N}: f(z)=a\}\right) $, then we say  $f$, $g$ share the set $S$ CM\\(IM) if $E_{f}(S)=E_{g}(S)$ $\left(\ol  E_{f}(S)=\ol E_{g}(S)\right) $. %\\Evidently, if $S$ contains only one element, then it coincides with the usual definition of CM(IM)  sharing of values. 
	\end{defi}
	\begin{defi}\cite{lahiriNagoya, Lahiri CVEE}
		Let k be a non-negative integer or infinity. For $a\in\overline{\mathbb{C}}$  we
		denote by $E_k(a; f)$ the set of all a-points of f, where an $a$-point of multiplicity $m$ is counted
		$m$ times if $m\leq k$ and $k + 1$ times if $m > k$. If $E_k(a; f) = E_k(a; g)$, we say that $f$, $g$ share
		the value a with weight $k$.
	\end{defi}
	We write $f$, $g$ share $(a,k)$ to mean that $f$, $g$ share the value a with weight $k$. Clearly if
	$f$, $g$ share $(a,k)$ then $f$, $g$ share $(a,p)$ for any integer $p$, $0\leq p < k$. Also we note that $f$, $g$
	share a value $a$ IM or CM if and only if $f$, $g$ share $(a,0)$ or $(a,\infty)$ respectively.
	\begin{defi}\cite{lahiriNagoya}
		For $S\subset\mathbb{\ol C}$ we define $E_f(S,k)=\cup_{a\in S}E_k(a;f)$, where k is a non-negative integer $a\in S$ or infinity. Clearly $E_f(S)=E_f(S,\infty)$ and $\overline{E}_f(S)=E_f(S,0)$. If $E_f(S,k)=E_g(S,k)$, then we say that $f$ and $g$ share the set $S$ with weight $k$.
	\end{defi}
	Obviously Definition 1.3 and Definition 1.4 are the refined notions of Definition 1.1 and Definition 1.2 respectively. However, now we go for the results as mentioned above. %However, now we recall the first result in this direction due to Steuding. 
	%\begin{theoA}\cite{steuding2007value}
	% 	If two L-functions     	$\mathcal{L}_1$ and $\mathcal{L}_2$ with $a(1)=1$  share a complex value $c\ne \infty$ CM, then $\mathcal{L}_1= \mathcal{L}_2$.
	%\end{theoA}
	%Since every L-function have meromorphic continuation in $\mathbb{C}$, so natural quest for the uniqueness of a meromorphic function and an L-function enters into the course of uniqueness theory vis-a-vis value distribution theory. Since an L-function can have at most one pole in $\mathbb{C}$, so it is reasonable to study the uniqueness of L-functions with meromorphic functions having finitely many poles. Pertinent to that, in 2010 Li  proved the following uniqueness theorem.
	\begin{theoA}\cite{li2010}
		Let a and b be two distinct
		finite values, and let f be a meromorphic function in
		the complex plane such that f has finitely many poles
		in the complex plane. If f and a non-constant L-function
		$\mathcal{L}$ share $(a,\infty)$ and $(b,0)$, then $\mathcal{L} = f$.
	\end{theoA}
	%\begin{rem}
	%	The number two of the finite shared values in Theorem B is the best possible, as can be easily verified with the help of the following two functions. \\Let $\mathcal{L} = \zeta$ and $f = \zeta e^g,$ where $g$ is an entire function. Observe that $\mathcal{L}$ and $f$ share $0$ CM but $\mathcal{L}\neq f$.
	%\end{rem}
	%Later on in 2013, Li-Yi proved the following uniqueness result considering the three value sharing.
	%\begin{theoC}\cite{li-yi}
	%	Let $a$, $b$ and $c$ be three distinct finite values, and let f be a transcendental meromorphic function in	$\mathbb{C}$ having finitely many poles in $\mathbb{C}$. If f and a non-constant L-function	$\mathcal{L}$ share $a, b, c$ IM, then $\mathcal{L} = f$.\end{theoC}
In 2018, taking the famous Gross Problem \cite{fg} into account,
Yuan, Li and Yi \cite{yuan2018} generalized the uniqueness problem of $\mathcal{L}$ and $f$ under the aegis of shared sets and   proposed the following inevitable question.
\begin{ques}\label{q1}\cite{yuan2018}
What can be said about the relationship between a meromorphic function $f$ and an L-function $\mathcal{L}$ if $f$ and $\mathcal{L}$ share one or two sets?
\end{ques}
Apropos of {\em Question \ref{q1}}, Yuan, Li and Yi provided the following result.
\begin{theoB}\cite{yuan2018}
Let $Q(z) = z^n + az^m + b$, where $a$, $b$ are non-zero constants with $\gcd(m,n)=1$ and $n \geq 2m + 5$. Further suppose  $f$ be a non-constant meromorphic function having finitely many poles and $\mathcal{L}$
be a non-constant L-function such that $E_{f}(S,\infty)=E_{\mathcal{L}}(S,\infty)$, where $S=\{z:Q(z)=0\}$. Then $f = \mathcal{L}$.
\end{theoB}
Later on,  Sahoo-Sarkar \cite{ss} improved  {\em Theorem B} by relaxing the weight  of the shared set from CM to weight 2.
\begin{theoC}\cite{ss}
Let $S$ be defined same as in {\em Theorem B} and $n \geq 2m+5$. Suppose $f$ be a non-constant meromorphic
function having finitely many poles in $\mathbb{C}$ and  $\mathcal{L}$ be a non-constant L-function. If $f$ and $\mathcal{L}$
share $(S, 2)$, then $f = \mathcal{L}$.
\end{theoC}
At the same time, Sahoo-Halder also considered the same problem with weight `$0$' and proved the following result. 
\begin{theoD}\cite{sh}
Let $S$ be defined same as in {\em Theorem B} and $n \geq \max\{2m+5,4q+9\},$ where $q=n-m\geq1$. Let $f$ be a non-constant meromorphic
function having finitely many poles in $\mathbb{C}$ and  $\mathcal{L}$ be a non-constant L-function. If $f$ and $\mathcal{L}$
share $(S, 0)$, then $f = \mathcal{L}$.
\end{theoD}
Though  {\em Theorem C} and {\em Theorem D} were later proved to be wrong by Banerjee-Kundu in \cite{bk}. Hence Banerjee-Kundu provided the following result rectifying the gaps of  these theorems.
\begin{theoE}\cite{bk}
Let $S$ be defined as in {\em Theorem B}, f be a non-constant meromorphic function having finitely
many poles in $\mathbb{C}$ and $\mathcal{L}$ be a non-constant L-function such that $E_f(S, t) = E_\mathcal{L}(S, t)$. If
\begin{enumerate}
	\item[(i)] $t \geq 2$ and $n \geq 2m + 5$, or
	\item[(ii)] $t = 1$ and $n \geq 2m + 6$, or
	\item[(iii)] $t=0$ and $n\geq 2m+11$, \\then $ f = \mathcal{L}$.
\end{enumerate} 	
\end{theoE}
In the same paper, Banerjee-Kundu also proved the following result analogous to {\em Theorem E}.
\begin{theoF}\cite{bk}
Let $S=\{z:z^{n}+az^{n-m}+b=0\}$, where $a,b$ are non-zero constants and $\gcd(n,m)=1$.  Let $f$ be a non-constant meromorphic function having finitely
many poles in $\mathbb{C}$ and $\mathcal{L}$ be a non-constant L-function such that $E_f(S, t) = E_\mathcal{L}(S, t)$. If
\begin{enumerate}
	\item[(i)] $t \geq 2$ and $n \geq 2m + 5$, or
	\item[(ii)] $t = 1$ and $n \geq 2m + 6$, or
	\item[(iii)] $t=0$ and $n\geq 2m+11$, \\then $ f = \mathcal{L}$.
\end{enumerate}   	\end{theoF}
Now from {\em Theorem B-F} one would naturally have the following observations.
\begin{enumerate}
\item[(i)] All the authors always used one of the following polynomials. \be\label{e000} P(z)=z^{n}+az^{m}+b\;\; or\;\; P(z)=z^{n}+az^{n-m}+b,\ee where $a,b$ are non-zero constants and $\gcd(n,m)=1$.  
\item[(ii)] The authors always considered a special class  of meromorphic functions; i.e., meromorphic functions having finitely many poles for the uniqueness of the same with an L-function.
\item[(iii)] The authors also considered that the set of zeros of the polynomials given by (\ref{e000}) may have multiple zeros. 
\end{enumerate}
%     Note that the set $S$  used in {\em Theorem C-G}  are generated from the zeros of the polynomial \be\label{e000} P(z)=z^{n}+az^{m}+b\;\; or\;\; P(z)=z^{n}+az^{n-m}+b,\ee where $a,b$ are non-zero constants and $\gcd(n,m)=1$. In \cite[see Lemma 4]{bk}, authors proved that these polynomials are critically injective and they may have  multiple zero but that must be one in number. On this occasion let us  invoke the definition of critically injective polynomial.
\vspace{0.1in}\par  %Observe that the following points come out of the above discussions.
%    \begin{enumerate}
%    	\item[(i)] All the authors always used one of the polynomials given by (\ref{e000}).  
%    	\item[(ii)] The authors always used the set of zeros of critically injective polynomials to show  the  uniqueness of $f$ and $\mathcal{L}$.
%\item[(iii)] In the above theorems authors have improved the previous results by relaxing the nature of sharing of the sets. % not by reducing the cardinalities of the sets.
%\item[(iv)] The authors also considered the set of zeros of the polynomials having multiple zeros. 
%  \end{enumerate}
Apropos of observation (i), One would naturally raise the following question.
\begin{ques}
Are these the only polynomials  given by (\ref{e000}) whose set of zeros provide uniqueness of $f$ and $\mathcal{L}$?%	Are these the only sets generated by either of the polynomials given by (\ref{e000}) for the uniqueness of $f$ and $\mathcal{L}$? 
\end{ques} 
If the answer is no. Then we would be happy to have the answer of the following general question.
\begin{ques}
Can we have a general characterization of polynomials  whose set of zeros provide uniqueness of $f$ and $\mathcal{L}$?
\end{ques}  
With respect to observation $(ii)$, we would like to recall that {\em Question \ref{q1}} was made for general meromorphic functions, whereas all the theorems (Theorem B-F) under the same question were treated only for meromorphic functions having finitely many poles. In this situation,  one would naturally have the following query.
\begin{ques}\label{q12}
Is it possible to extend observation $(ii)$ for the whole class of meromorphic functions in order to have a full answer of {\em Question \ref{q1}}?
\end{ques}

%     \begin{ques}\label{q5}
%     Does there exist any non-critically injective polynomial whose set of zeros provide the uniqueness of $f$ and $\mathcal{L}$?
%    \end{ques}
%    Pertinent to observation (iii) the following  questions become inevitable.
%    \begin{ques}\label{q3}
%   	Can we have the answer of {\em Question \ref{q1}} under more relaxed sharing hypothesis than that obtained in the latest results {\em Theorem G-H}?
%  \end{ques}
% \begin{ques}\label{q6}
%Can we have  a set with lesser cardinality than that obtained in the latest results {\em Theorem G-H} for the uniqueness of $f$ and $\mathcal{L}$?	
%\end{ques}
Finally with respect to observation (iii), we have the following remark.
\begin{rem}
	Very recently	in \cite[see paragraph between Theorem H and Theorem I]{bk1} Banerjee-Kundu have shed a light on  the fact that all the results %obtained till date in this direction of shared set problems for the uniqueness of $f$ and $\mathcal{L}$; i.e.,
	from {\em Theorem B-F} are valid only when the polynomials  have simple zeros or the polynomials have multiple zeros with weight of the shared set as $0$ and not valid otherwise. Moreover, a close look in the proof of these theorems reveals that the concept of multiple zeros of polynomials does not provide us with any such analytical or technical advantage. %Though {\em Theorem F} has a different flaw contradicting their own conclusion of cardinality $n\geq\max\{2m+5,4q+9\}$ which is analysed in \cite[Remark 3]{bk}.     	 \\\\have an analytical gap while considering multiple zero of the polynomials and the sharing of the sets with some non-zero weight. Thus conclusion of {\em Theorem D-H} (except Theorem F) become false when the multiple zero of the generating polynomials are taken into account and the sharing of the sets with some non-zero weight. But in the same scenario, the results obtained with IM sharing of the sets are correct; i.e., conclusion (iii) of {\em Theorem G-H} and {\em Theorem F}. Though {\em Theorem F} has a different flaw contradicting their own conclusion of cardinality $n\geq\max\{2m+5,4q+9\}$ which is analysed in \cite[Remark 3]{bk}. Another point is that all these results are true when the polynomial has only simple zeros. 
\end{rem}

Hence in this paper, we  solely concentrate on the polynomials having only simple zeros and answer all the above questions from {\em Question \ref{q1}-\ref{q12}} affirmatively. % which improve all the existing results from {\em Theorem D-H}.
In fact, we present   general criterions for any general polynomial so that the set of zeros of the same would provide the uniqueness of  $\mathcal{L}$ with a general meromorphic function $f$ instead of a meromorphic function having finitely many poles. That is, we extend all these results to the most general setting. As a consequence, all the above results comes under the special cases of our main theorems. %in view of generalising the special class of meromorphic function to the whole class of meromorphic function as well as finding general criterions for any general polynomial for the uniqueness of  $\mathcal{L}$ with a general meromorphic function $f$. 
Moreover, our main theorems significantly improve all the above results obtained for ignoring multiplicities providing better estimates of least cardinalities; i.e., $11$ instead of $13$.   In a nutshell, our main theorems bring all the existing theorems under a single umbrella in a more improved version with extent to the most general setting. 
%   \begin{defi}
	%   	Let $P(z)$ be a polynomial such that $P^{'}(z)$ has mutually $r$ distinct  zeros given by $d_{1}, d_{2}, \ldots, d_{r}$ with multiplicities $q_{1}, q_{2}, \ldots, q_{r}$ respectively. Then $P(z)$ is said to be a critically injective polynomial if $P(d_i)\not =P(d_j)$ for $i\not=j$, where $i,j\in \{1,2,\cdot\cdot\cdot,r\}$.
	%  \end{defi}
% Any polynomial which is not critically injective is called a non-critically injective polynomial.

\vspace{0.1in}\par In the {\bf ``Application"} section we have  exhibited a number of examples proving all our claims to be true and showing the far reaching applications of our  results.%We also conclude that the answer of {\em Question 1.5} can not be obtained for any polynomial in all the existing methods used to prove all the above results. 	
\par Before going to our main results, we make a short discussion on the structure of a general
polynomial as this will play a crucial role throughout the paper. 
%Let us now place an example to clarify the above discussion.
\vspace{0.1in}\par Let us consider the following general polynomial $P(z)$ of degree $n$ having only simple zeros.  	\bea\label{el0} P(z)=a_nz^n+a_{n-1}z^{n-1}+\ldots+a_1z+a_0,\eea
where $ a_0,a_1,\dots,a_n$ are  complex numbers with $a_n,a_0\neq 0$, $a_{i}$ being the first non-vanishing coefficient  from $a_{n-1}, a_{n-2},\ldots,a_{1}$. Let \be\label{el00}S=\{z:P(z)=0\}.\ee
Observe that (\ref{el0}) can be written in the form \be\label{el000}P(z)=a_{n}\prod\limits_{i=1}^{p}(z-\alpha_{i})^{m_{i}}+a_{0},\ee where $p$ denotes the   number of distinct zeros of $P(z)-a_{0}$. Let us also denote by $s$ the number of distinct zeros of $P^{'}(z)$. Hence we would have 
\be\label{el0000}P^{'}(z)=na_{n}\prod\limits_{i=1}^{s}(z-\eta_{i})^{r_{i}},\ee where $r_{i}$ denotes the multiplicities of distinct zeros of $P^{'}(z)$.
\vspace{0.2in}\par 	
Set 
\bea\label{el2} R(z)&=& -\frac{a_nz^n}{a_iz^i+a_{i-1}z^{i-1}+\ldots+a_1z+a_0}\\\nonumber&=&- \frac{a_{n}z^n}{a_{i}\prod\limits_{j=1}^{k}(z-\beta_j)^{m_{j}}}\\\nonumber&=&\d-\frac{a_{n}z^{n}}{\phi(z)},\eea
where $a_{0},a_{1},\ldots,a_{n}, a_{i}$ are as defined in (\ref{el0}) and $\beta_1,\beta_2,\ldots,\beta_k$ are the roots of the equation $$\phi(z)=a_iz^i+a_{i-1}z^{i-1}+\ldots+a_1z+a_0=0,$$ with multiplicities $m_1,m_2,\ldots,m_k$. Clearly \be\nonumber R(z)-1=\d-\frac{P(z)}{\phi(z)},\ee where $P(z)$ is defined by (\ref{el0}) and obviously $P(z)$ and $\phi(z)$ do not share any common zero. Hence $S$ as defined in (\ref{el00}) can be treated as \be\label{el02}S=\{z:P(z)=0\}=\{z:R(z)-1=0\}.\ee  Let $R^{'}(z)$ has $l$ distinct zeros say $\delta_{1},\delta_{2},\ldots,\ldots,
\delta_{l}$ with multiplicities $q_{1},q_{2},\ldots,q_{l}$ respectively. Then  
From (\ref{el2}) we would have 
\bea \label{el3} R'(z)&=&\frac{\gamma\prod\limits_{j=1}^{l}(z-\delta_j)^{q_{j}}}{\prod\limits_{j=1}^{k}(z-\beta_j)^{p_{j}}},\eea where $\gamma\in\mathbb{C}-\{0\}$ and $p_{j}\in\mathbb{N}$ for all $j\in\{1,2,\ldots,k\}$. %$p_{j}>m_{j}$ for $j\in\{1,2,\ldots,k\}$.
\begin{rem}
	Observe that in the definition (\ref{el0}) of the general polynomial $P(z)$, the condition $a_{i}\neq 0$ for at least one $i\in \{1,2,\ldots,n-1\}$ is necessary. Because otherwise we would find a non-constant L-function $\mathcal{L}$ and a non-constant meromorphic function $f$ which share the set $S=\{z:P(z)=0\}$ CM but $f\neq \mathcal{L}$.
	\par  For example, let  $a_{i}=0$ for all $i\in \{1,2,\ldots,n-1\}$. Then $S=\{z: a_{n}z^{n}+a_{0}=0\}$. Consider a non-constant L-function $\mathcal{L}$ and  a non-constant meromorphic function $f$ such that $f=\zeta\mathcal{L}$, where $\zeta$ is the nth root of unity such that $\zeta \ne 1 $. Then clearly, $$a_{n}f^{n}+a_{0}=a_{n}\mathcal{L}^{n}+a_{0};$$ $$i.e.,\;\; \prod\limits_{i=1}^{n}(f-\sigma_{i})=\prod\limits_{i=1}^{n}(\mathcal{L}-\sigma_{i}),$$ where $\sigma_{i}\in S$ for $i=\{1,2,\ldots,n\}$. That is $f$ and $\mathcal{L}$ share $S$ CM but $f\ne \mathcal{L}$.
\end{rem}  

\section{Main Results}   	
The following two theorems are the main results of this paper.
% first main result of this paper as the following theorem.
\begin{theo}\label{t1}
	Let $P(z)$ be given by (\ref{el000}) with $p\geq2$, and $S$,$s$ be defined by (\ref{el00}),(\ref{el0000}) respectively.
	%\begin{enumerate}
	%   \item[(a)] $p\geq 3$, or
	%  \item [(b)] $p=2$ and $gcd(m_i,n)=1$ for at least one of the $m_i$'s such that $n=\sum\limits_{i=1}^{2} m_i \geq 3$, or
	% \item[(c)] $p=2$ and $gcd(m_i,n)\ne 1$ for each $m_i$ such that $n\geq (b_1+b_2)+1$, where $b_i=gcd(m_i,n)$,
	%\end{enumerate}
	Suppose $f$ be a non-constant meromorphic function and  $\mathcal{L}$
	be a non-constant L-function sharing $(S,t)$.
	% such that $E_{f}(S, 2)=E_{\mathcal{L}}(S,2)$.
	Then for $n \geq \max \left\{2p+2,2s+5\right\}$, when $t\geq2$; $n \geq \max \left\{2p+2,2s+6\right\}$, when $t=1$ and for\\ $n \geq \max \left\{2p+2,2s+11\right\}$, when $t=0$; the following are equivalent :
	\begin{itemize}
		\item[(i)] $P(f)=P(\mathcal{L}) \implies f=\mathcal{L}$ ;
		\item[(ii)] $E_f(S, t)=E_\mathcal{L}(S,t)\implies f=\mathcal{L}$.
		%		\item[(iii)] $P(f)=cP(\mathcal{L}) \implies f=\mathcal{L}$,  where $c$ is a non-zero complex number. 
	\end{itemize}
\end{theo}
%Following theorem is the second main result of the paper.
\begin{theo}\label{t2}
	Let $R(z)$ be defined by (\ref{el2}) with $k\in \mathbb{N}$ and $S, l$ be defined by (\ref{el02}), (\ref{el3}) respectively with \begin{enumerate}
		\item[(a)] $k\geq2$, or
		\item[(b)] $k=1$ and  $n\geq 3b_1+1$, where $b_1=gcd(m_1,n)$.
	\end{enumerate} 
	Let $f$ be a non-constant meromorphic function and  $\mathcal{L}$
	be a non-constant L-function sharing $(S,t)$.
	Then for $n \geq \max \left\{2k+4,2l+5\right\}$, when $t\geq2$;\\ $n \geq \max \left\{2k+4,2l+6\right\}$, when $t=1$ and for $n \geq \max \left\{2k+4,2l+11\right\}$, when $t=0$; the following are equivalent :
	\begin{itemize}
		\item[(i)] $R(f)=R(\mathcal{L}) \implies f=\mathcal{L}$ ;
		\item[(ii)] $E_f(S, t)=E_\mathcal{L}(S,t)\implies f=\mathcal{L}$.	\end{itemize}	
	
\end{theo} 
Next we state the following two corollaries with respect to {\em Theorem \ref{t1}} and {\em Theorem \ref{t2}} respectively for meromorphic functions having finitely many poles. %In fact, these corollaries, in the application part would confirm us that {\em Theorem B-F} are nothing but the special cases of these corollaries with significant improvement with respect to weight $0$.   
\begin{cor}\label{c1}
	Let $P(z)$, $p$, $S$, $s$ be as defined in {\em Theorem \ref{t1}}.
	%\begin{enumerate}
	%   \item[(a)] $p\geq 3$, or
	%  \item [(b)] $p=2$ and $gcd(m_i,n)=1$ for at least one of the $m_i$'s such that $n=\sum\limits_{i=1}^{2} m_i \geq 3$, or
	% \item[(c)] $p=2$ and $gcd(m_i,n)\ne 1$ for each $m_i$ such that $n\geq (b_1+b_2)+1$, where $b_i=gcd(m_i,n)$,
	%\end{enumerate}
	Suppose $f$ be a non-constant meromorphic function having finitely many poles and  $\mathcal{L}$
	be a non-constant L-function sharing $(S,t)$.
	% such that $E_{f}(S, 2)=E_{\mathcal{L}}(S,2)$.
	Then for $n \geq \max \left\{2p+1,2s+3\right\}$, when $t\geq2$; $n \geq \max \left\{2p+1,2s+4\right\}$, when $t=1$ and for $n \geq \max \left\{2p+1,2s+7\right\}$, when $t=0$; the following are equivalent :
	\begin{itemize}
		\item[(i)] $P(f)=P(\mathcal{L}) \implies f=\mathcal{L}$ ;
		\item[(ii)] $E_f(S, t)=E_\mathcal{L}(S,t)\implies f=\mathcal{L}$.
		%		\item[(iii)] $P(f)=cP(\mathcal{L}) \implies f=\mathcal{L}$,  where $c$ is a non-zero complex number. 
	\end{itemize}
\end{cor}
%Following theorem is the second main result of the paper.
\begin{cor}\label{c2}
	Let $R(z)$, $k$,  $S, l$ be as defined in {\em Theorem \ref{t2}} with \begin{enumerate}
		\item[(a)] $k\geq2$, or
		\item[(b)] $k=1$ and  $n\geq 2b_1+1$, where $b_1=gcd(m_1,n)$.
	\end{enumerate} 
	Let $f$ be a non-constant meromorphic function having finitely many poles and  $\mathcal{L}$
	be a non-constant L-function sharing $(S,t)$.
	Then for $n \geq \max \left\{2k+3,2l+3\right\}$, when $t\geq2$; $n \geq \max \left\{2k+3,2l+4\right\}$, when $t=1$ and for $n \geq \max \left\{2k+3,2l+7\right\}$, when $t=0$; the following are equivalent :
	\begin{itemize}
		\item[(i)] $R(f)=R(\mathcal{L}) \implies f=\mathcal{L}$ ;
		\item[(ii)] $E_f(S, t)=E_\mathcal{L}(S,t)\implies f=\mathcal{L}$.	\end{itemize}	
	
\end{cor}   
\section{Applications}
In this section, we exhibit different examples applying our main theorems and corollaries which clearly shows the broad perspectives of our results with special attention to the fact that the results  {\em Theorem B-F} comes under the special cases of our results. In addition, our results significantly improve {\em Theorem B-F} in case of weight `$0$' of the shared set. Below we explain this fact via the following two examples.
%   \\\\generalize the existing results in such a way that all the variants of polynomials can be accumulated  under single shed.
\begin{exm}\label{ex2}
	Consider the polynomial 
	\be\nonumber
	P(z)=z^n+az^{n-m}+b= z^{n-m}(z^{m}+a)+b,
	\ee
	where $n, m$  are relatively prime inegers with $n-m>1$and a, b are non-zero constants such that the polynomial has no multiple zero. Suppose $S=\{z:P(z)=0\}$.
	Here $$p=m+1\geq 2 \;\ \text{and }\;\ gcd(n,n-m)=1.$$Also \be\nonumber
	P'(z)=z^{n-m-1}(nz^{m}+a(n-m));
	\ee
	i.e., $$s=m+1.$$
	
	
	Suppose that $P(f)=P(\mathcal{L})$, for any non-constant meromorphic function f and a non-constant L-function $\mathcal{L}$,  then we have
	\be\label{411}
	f^n-\mathcal{L}^n=-a(f^{n-m}-\mathcal{L}^{n-m}).
	\ee
	If $f^n \not\equiv \mathcal{L}^n$, then we can rewrite (\ref{411}) as
	\be\label{412}
	\mathcal{L}^{m}=-a\frac{(h-v)(h-v^2)...(h-v^{{n-m}-1})}{(h-u)(h-u^2)...(h-u^{n-1})},
	\ee
	where $h =\displaystyle\frac{f}{\mathcal{L}}$, $u = exp(2\pi i/n)$ and $v = exp(2\pi i/(n-m))$. Noting that n and (n-m) are relatively prime positive
	integers, then the numerator and denominator of (\ref{412}) have no common factors. Since $\mathcal{L}$ has atmost one pole at $z = 1$ in
	the complex plane, and whenever $n\geq 5$ we can see that there exists at least three distinct roots of $h^n = 1$ such that they are Picard
	exceptional values of h$,$ and so it follows by (\ref{412}) that h and thus $\mathcal{L}$ are constants, which is impossible.
	
	Therefore, we must have $f^n = \mathcal{L}^n$. Then by (\ref{411}) we also have $f^{n-m} = \mathcal{L}^{n-m}$. Since n and (n-m) are relatively prime positive integers, we deduce that $f = \mathcal{L}$. Thus we see that $P(f)=P(\mathcal{L})\implies f=\mathcal{L}$, when $n\geq 5$. 
	
	Now we apply {\em\bf Theorem \ref{t1}} to find the minimum value of $n$ for which we can say that  $E_f(S, t)=E_\mathcal{L}(S,t)\implies f=\mathcal{L}$.
	
	Therefore,  $E_f(S, t)=E_\mathcal{L}(S,t)\implies f=\mathcal{L}$ for  
	\begin{enumerate}
		\item $n \geq \max\{ 2m+4,2m+7\}=2m+7$ when $t\geq 2 $,
		\item  $n \geq \max\{ 2m+4,2m+8\}=2m+8$ when $t=1 $, and
		\item  $n \geq \max\{ 2m+4,2m+13\}=2m+13$ when $t=0 $.
	\end{enumerate}
	
	Again if $f$ is a non-constant meromorphic function with finitely many poles and $\mathcal{L}$ is a non-constant $L$ function, then by {\em\bf Corollary \ref{c1}} we have $E_f(S, t)=E_\mathcal{L}(S,t)\implies f=\mathcal{L}$ for  
	\begin{enumerate}
		\item $n \geq \max\{ 2m+3,2m+5\}=2m+5$ when $t\geq 2 $,
		\item  $n \geq \max\{ 2m+3,2m+6\}=2m+6$ when $t=1 $, and
		\item  $n \geq \max\{ 2m+3,2m+9\}=2m+9$ when $t=0 $.
	\end{enumerate} 
\end{exm}
\begin{exm}\label{ex5}
	Consider the polynomial \bea\nonumber  P(z) = z^n + az^m + b,\eea where $m$ and $n$ are positive integers such that $n \geq m+4$, $a$ and $b$ are finite non-zero complex numbers with $\frac{b^{n-m}}{a^m}\ne \frac{(-1)^nm^m(n-m)^{n-m}}{n^n}$. Then $P(z)$ has only simple zeros. Let $S$ denotes the set of zeros of $P(z)$. Suppose \bea\label{eyi} 
	R(z)=- \frac{z^n}{az^m+b}.\eea
	Then we find that $S=\{z:R(z)-1=0\}$.
	From (\ref{eyi}) we have \bea \nonumber
	R'(z)=- \frac{z^{n-1}[a(n-m)z^m+bn]}{(az^m+b)^2}.\eea 
	Now for a non-constant meromorphic function $f$ and a non-constant L-function $\mathcal{L}$ consider $R(f)=R(\mathcal{L})$. Then we have  \bea\label{e49} 
	\frac{f^n}{af^m+b}&=&\frac{\mathcal{L}^n}{a\mathcal{L}^m+b}\nonumber\\
	\implies a(f^n\mathcal{L}^m-f^m\mathcal{L}^n)-b(\mathcal{L}^n-f^n)&=&0.\eea
	Let   $h=\frac{f}{\mathcal{L}}$. Suppose that $h $ is a non-constant meromorphic function. Then from (\ref{e49}) we have
	\bea \label{e111}&&
	ah^{m}\mathcal{L}^{n+m}(h^{n-m}-1)+b\mathcal{L}^n(h^n-1)=0\\\nonumber&&
	\implies \mathcal{L}^m=-\frac{b(h^n-1)}{ah^{m}(h^{n-m}-1)}=-\frac{b}{a}\frac{(h-u)(h-u^2)...(h-u^{n-1})}{h^m(h-v)(h-v^2)...(h-v^{n-m-1})},
	\eea
	where  $u = exp(2\pi i/n)$, and $v = exp(2\pi i/(n-m))$. Since $n$ and $m$ are co-prime, so is $n$ and $n-m$. Hence the numerator and denominator of (\ref{e111}) have no common factors. Further, the function $\mathcal{L}$ has atmost one pole  in
	the complex plane, it follows that $h$ has atleast $n-m-1$ picard exceptional values among $\{0,v,v^2,...,v^{n-m-1}\}$.
	
	Clearly this is a contradiction as $n\geq m+4$. Hence  $h$ is  constant. Thus from (\ref{e111}) we must have  $h^{n}=1=h^{n-m}$ which in turn implies $h=1$; i.e., $f=\mathcal{L}$. \vspace{0.1in}\par Now we count the cardinality of the set $S$ for which  $E_f(S, t)=E_\mathcal{L}(S,t)\implies f=\mathcal{L}$. In this case, for $R(z)$ we have  $$l=m+1,\;\; k=m.$$  Therefore we obtain that $R(z)$ satisfies condition (a)[when $m\geq2$], or (b)[when m=1] and (i)  of
	{\em\bf Theorem \ref{t2}}. Hence  $E_f(S, t)=E_\mathcal{L}(S,t)\implies f=\mathcal{L}$ for \begin{enumerate}
		\item $n\geq \max\{2m+4,2m+7\}=2m+7$ when $t\geq 2$ ,
		\item $n\geq \max\{2m+4,2m+8\}=2m+8$ when $t= 1$, and 
		\item $n\geq \max\{2m+4,2m+13\}=2m+13$ when $t=0$.
	\end{enumerate}
	Now let $f$ be a non-constant meromorphic function with finitely many poles and $\mathcal{L}$ be a non-constant $L$ function. Then by {\em\bf Corollary \ref{c2}} we have $E_f(S, t)=E_\mathcal{L}(S,t)\implies f=\mathcal{L}$ for  
	\begin{enumerate}
		\item $n \geq \max\{ 2m+3,2m+5\}=2m+5$ when $t\geq 2 $,
		\item  $n \geq \max\{ 2m+3,2m+6\}=2m+6$ when $t=1 $, and
		\item  $n \geq \max\{ 2m+3,2m+9\}=2m+9$ when $t=0 $.
	\end{enumerate} 		
\end{exm} 	
\begin{rem}
	Observe that the conclusions obtained in {\em Example \ref{ex2}} and {\em Example \ref{ex5}} for meromorphic functions having finitely many poles are at par with all the existing results from {\em Theorem B-F} with an improvement in the case of weight ` $0$' $($ i.e., the cardinality is reduced from $2m+11$ to $2m+9)$.
\end{rem}
Another point we would like to explore on this occasion that the polynomials used in {\em Theorem B-F} are all critically injective \cite[see Lemma 4]{bk}. For the sake of our readers, we recall the definition of this particular type of polynomials. 
\begin{defi}
	Let $P(z)$ be a polynomial such that $P^{\prime}(z)$ has mutually $r$ distinct  zeros given by $d_{1}, d_{2}, \ldots, d_{r}$ with multiplicities $q_{1}, q_{2}, \ldots, q_{r}$ respectively. Then $P(z)$ is said to be a critically injective polynomial if $P(d_i)\not =P(d_j)$ for $i\not=j$, where $i,j\in \{1,2,\cdot\cdot\cdot,r\}$.
	\par 	Any polynomial which is not critically injective is called a non-critically injective polynomial.
\end{defi}
Obviously, any polynomial is either critically injective or non-critically injective. So, taking this fact into account, we have framed our theorems in such a way that the non-critically injective polynomials can also be fit into our theorems.  Even, the polynomials which are still unknown to be critically injective or non-critically injective also fit into our theorems. That is our results count all kinds of polynomials answering {\em Question \ref{q1}-\ref{q12}} affirmatively. Below we provide examples of each these three types of polynomials respectively, justifying all the above claims.	
\begin{exm}
	Let us consider the  polynomial :
	\bea\label{e44.0} P(z)=z^n+az^{n-m}+bz^{n-2m}+c,\eea
	where $a, b, c\in\mathbb{C}^*$ be such that $P(z)$ has no multiple root, $gcd(m, n) = 1$ with $n\geq2m+4$ and $\frac{a^2}{4b}=\frac{n(n-2m)}{(n-m)^2}$, $c\ne \frac{\beta_i\beta_j}{\beta_i+\beta_j}$. Here $\beta_i=-(c_i^n+ac_i^{n-m}+bc_i^{n-2m}),$ where $c_i$ are the roots
	of the equation $nz^{2m} + (n- m)az^m + b(n-2m) = 0$, for $i = 1, 2, \ldots, 2m$. 
	Suppose $S$ denotes the set of zeros of (\ref{e44.0}).
	\vspace{0.1in}\par  Obviously,  $P(z)$ has only simple zeroes and it is critically injective \cite[see Lemma 2.7]{3a}. 
	From (\ref{e44.0}) we have \bea\label{e440} P'(z)&=&z^{n-2m-1}[nz^{2m}+a(n-m)z^m+b(n-2m)]\\\nonumber&&=nz^{n-2m-1}\left(z^m+\frac{a(n-m)}{2n}\right)^2.\eea
	From (\ref{e44.0}) and (\ref{e440}) we find that $$p=2m+1 \;\;\; \mbox{and} \;\;\;s=m+1.$$
	In \cite[see proof of Theorem 1.1]{3a}, it is also proved that $P(f)=P(g)$ implies $f=g$ for $n\geq 2m+4$, where $f$and $g$ are non-constant meromorphic functions. Hence for a non-constant meromorphic function $f$ and an L-function  $\mathcal{L}$ we have $P(f)=P(\mathcal{L})\implies f=\mathcal{L}$ when $n\geq 2m+4$. Thus $P(z)$ satisfies the condition (i) of {\em\bf Theorem \ref{t1}} of the present paper and  hence $E_f(S, t)=E_{\mathcal{L}}(S,t)\implies f=\mathcal{L}$ for
	\begin{enumerate}
		\item $n\geq\max\{4m+4,2m+7\}\geq 9$ when $t \geq 2$,
		\item  $n\geq\max\{4m+4,2m+8\}\geq 10$ when $t = 1$, and
		\item  $n\geq\max\{4m+4,2m+13\}\geq 15$ when $t = 0$.
	\end{enumerate}
	Again if $f$ is a non-constant meromorphic function with finitely many poles and $\mathcal{L}$ be a non-constant $L$ function, then by {\em\bf Corollary \ref{c1}} we have $E_f(S, t)=E_\mathcal{L}(S,t)\implies f=\mathcal{L}$ for  
	\begin{enumerate}
		\item $n \geq \max\{ 4m+3,2m+5\}=2m+5 \geq 7$ when $t\geq 2 $,
		\item  $n \geq \max\{ 4m+3,2m+6\}=2m+6 \geq 8$ when $t=1 $, and
		\item  $n \geq \max\{ 4m+3,2m+9\}=2m+9 \geq 11$ when $t=0 $.
	\end{enumerate} 
	
\end{exm}

\begin{rem}
	Note that the polynomial \be\label{efr}P(z)=\frac{(n-1)(n-2)}{2}z^n-n(n-2)z^{n-1}+\frac{n(n-1)}{2}z^{n-2}-c,\ee
	where $n\geq6,$ $c\neq0,1$, introduced by Frank-Reinders \cite{frank1998unique} comes as the special case of (\ref{e44.0}) for $m=1$, $a=\frac{-2n}{n-1}$, $b=\frac{n}{n-2}$ and $c\in\mathbb{C}-\{0,\frac{-1}{(n-1)(n-2)}\}$. Hence $E_f(S, t)=E_{\mathcal{L}}(S,t)\implies f=\mathcal{L}$ as $n\geq 9$ when $t \geq 2$, $n\geq 10$ when $t =1$ and $n\geq 15$ when $t = 0$, where $S$ denotes the set of zeros of (\ref{efr}) and $f$, $\mathcal{L}$ are non-constant meromorphic function  and a non-constant L-function respectively.
	%	\end{equation*} 	where $c\ne 0,1$ and $n\geq6$..
\end{rem}
%	In the next example we explore a non-critically injective polynomial in the direction of Theorem 	1.1.
\begin{exm}
Let

\begin{equation*}
	P(z) = z^n-4z^{n-1}+4z^{n-2}+c,
\end{equation*}
where $n(\geq 5)$ is odd,
$c\in \mathbb{C}$ such that $P(z)$ does not have any multiple zero. Also we have 
\bea\nonumber P'(z)=z^{n-3}\left(nz^2-4(n-1)z+4(n-2)\right).\eea
Here $P(z)$ is a non-critically injective polynomial \cite[see  (1.1)]{smtbilisi} and we see that $$p=2,\;\; s=3.$$ Suppose $S=\{z:P(z)=0\}$. 
Let $f$ and $\mathcal{L}$ be two non-constant meromorphic  and L-function respectively such that $$P(f) = P(\mathcal{L}).$$ Since $\mathcal{L}$ has at most one pole in $\mathbb{C}$, hence proceeding in the same line of proof as done in {\em Case-3 of Theorem 1.3} of \cite{smtbilisi} or {\em Example 4.4} of \cite{smfil21}  we also get here $f=\mathcal{L}$. 

Therefore $P(z)$ satisfies condition (i) of {\em\bf Theorem \ref{t1}}. Hence we conclude that $E_f(S, t)=E_\mathcal{L}(S,t)\implies f=\mathcal{L}$ when  
\begin{enumerate}
	\item $n\geq \max\{2.2+2,2.3+5\}=11$ for $t\geq 2$, 
	\item $n\geq \max\{2.2+2,2.3+6\}=12$ for $t=1$, and
	\item $n\geq \max\{2.2+2,2.3+11\}=17$ for $t=0$.
\end{enumerate}

Now let  $f$ be a non-constant meromorphic function with finitely many poles and $\mathcal{L}$ be a non-constant $L$ function. Then by {\em\bf Corollary \ref{c1}} we have $E_f(S, t)=E_\mathcal{L}(S,t)\implies f=\mathcal{L}$ for  
\begin{enumerate}
	\item $n \geq \max\{ 2.2+1,2.3+3\}=9$ when $t\geq 2 $,
	\item  $n \geq \max\{ 2.2+1,2.3+4\}=10$ when $t=1 $, and
	\item  $n \geq \max\{ 2.2+1,2.3+7\}=13$ when $t=0 $.
\end{enumerate} 
\end{exm}	 

\begin{exm}
Consider the polynomial 
\bea P(z)=az^n-n(n-1)z^2+2n(n-2)bz -(n-1)(n-2)b^2,\nonumber\eea
where $n(\geq6)$ is an integer and $a$, $b$ are two non-zero complex numbers satisfying
$ab^{n-2}\ne 1, 2$. In \cite[see section 5]{bl}, we can find that this polynomial is still uncertain to be critically injective or non-critically injective. Suppose $S=\{z:P(z)=0\}$.  It is obvious that $n(n-1)z^2-2n(n-2)bz +(n-1)(n-2)b^2=0;$ has two
distinct roots, say $\alpha_1$ and $\alpha_2$. 
Here  \bea\label{e41} R(z)=\frac{az^n}{n(n-1)(z-\alpha_1)(z-\alpha_2)}. \eea
Hence $S=\{z:R(z)-1=0\}$.
From (\ref{e41}) we have
\bea\nonumber R'(z)=\frac{(n-2)az^{n-1}(z-b)^2}{n(n-1)(z-\alpha_1)^2(z-\alpha_2)^2}.\eea
Let $f$  a non-constnat meromorphic function  and $\mathcal{L}$ be a non-constant L-function. Since every L-function is meromorphic in $\mathbb{C}$, so  $R(f)=R(\mathcal{L})\implies f=\mathcal{L}$ for $n\geq6$ directly follows from \cite[see page 67]{aljahari}.
We also find that in this case $$l=2,\;\;k=2. $$ Since $P(z)$ satisfies condition (a) and (i) of  {\em\bf Theorem \ref{t2}}. Hence we obtain $E_f(S, t)=E_\mathcal{L}(S,t)\implies f=\mathcal{L}$ for \begin{enumerate}
	\item $n\geq \max\{2k+4,2l+5\}=9$ when $t\geq2$,
	\item $n\geq \max\{2k+4,2l+6\}=10$ when $t=1$, and 
	\item for $n\geq \max\{2k+4,2l+11\}=15$ when $t=0$.
\end{enumerate}

Again if $f$ is a non-constant meromorphic function with finitely many poles and $\mathcal{L}$ is a non-constant L-function, then by {\em\bf Corollary \ref{c2}} we have $E_f(S, t)=E_\mathcal{L}(S,t)\implies f=\mathcal{L}$ for  
\begin{enumerate}
	\item $n \geq \max\{ 2k+3,2l+3\}=7$ when $t\geq 2 $,
	\item  $n \geq \max\{ 2k+3,2l+4\}=8$ when $t=1 $, and
	\item  $n \geq \max\{ 2k+3,2l+7\}=11$ when $t=0 $.
\end{enumerate} 
\end{exm}
%\begin{rem}\label{r2}
%	Obviously {\em Theorem G-H} are the latest results in this direction for simple zeros of the polynomials given by (\ref{e000}). In the application section ({\em Example \ref{ex5}} and {\em Example \ref{ex2}}), we shall show that the conclusions of {\em Theorem G-H} are true for $n\geq 7$ when $E_{f}(S,1)=E_{\mathcal{L}}(S,1)$, whereas the same is true in {\em Theorem G-H} for $n\geq 8$. Thus our result directly improves {\em Theorem G-H} by reducing the cardinality of the set $S$ when shared by the functions with weight $1$. We  also find that weight $2$ in {\em Theorem G-H} can be relaxed to weight $1$ keeping the carinality of the set fixed as an application of our result. Hence the answer of {\em Question \ref{q3}} is also obtained with improvement.  Moreover, in {\em Theorem G-H} the least cardinality of the sets when shared IM is 13 whereas the same result can be obtained when the cardinalities of the sets are 10, which is a significant improvement of {\em Theorem G-H}. Thus we obtain a threefold improvement of {\em Theorem G-H} by the application of our main results and obtain the answer of {\em Question \ref{q3}-\ref{q6}}. 
%	\vspace{0.1in}\par We shall also obtain similar results in the application section for other polynomials including critically injective polynomials, non-critically injective polynomials and even those polynomials which are still uncertain to be critically injective or non-critically injective (see {\em Example \ref{ex1}}, {\em Example \ref{ex3}} and {\em Example \ref{ex4}}). These results provide us the answers of {\em Question \ref{q4}-\ref{q5}} with improvements in the nature of sharing of the sets as well as the least cardinalities of the sets.
%  	\end{rem}
%  	\begin{rem}
% 		The reason behind proving two similar but different theorems are clarified in the first two paragraphs of section 5 named {\bf ``Conclusion and an Open Question"}.
%	\end{rem}

%\begin{defi}
%	Let $k$ be a positive integer or infinity. We denote by $N_k(r, b; f)$
%	the counting function of b-points of $f$, where an b-point of multiplicity $m$ is counted m
%	times if $m \leq k$ and $k$ times if $m > k$. That is\begin{equation*}
%	N_k(r, b; f) = \ol{N}(r, b; f) + \ol{N}(r, b; f |\geq 2) + \ldots + \ol{N}(r, b; f |\geq k).
%	\end{equation*}
%	\end{defi}
%\begin{defi}\cite{Lahiri CVEE}
%	We denote by $N_{2}(r,a;f)=\ol N(r,a;f)+\ol N(r,a;f|\geq2)$.
%\end{defi}  	
\section{Lemmas}
In this section we prove some lemmas which will be required for the proofs of our main results. Before stating these lemmas we shortly recall the following definitions for the convenience of our readers. %For standard definitions and notations we have already suggested our readers to follow \cite{wkh}. Furthermore, for the convenience of our readers, we explain the following notations which will be required in the proofs of the lemmas and main results. % for the sake of our convenience.
\begin{defi} \cite{15} Let $f$ and $g$ be two non-constant meromorphic functions such that $f$ and $g$ share $(1,0)$. Let $z_{0}$ be a $1$-point of $f$ with multiplicity $p$, a $1$-point of $g$ with multiplicity $q$. We denote by $\ol N_{L}(r,1;f)$ the reduced counting function of those $1$-points of $f$ and $g$ where $p>q$, by $N^{1)}_{E}(r,1;f)$ the counting function of those $1$-points of $f$ and $g$ where $p=q=1$. In the same way we can define $\ol N_{L}(r,1;g)$, $N^{1)}_{E}(r,1;g)$. In a similar manner we can define $\ol N_{L}(r,a;f)$ and $\ol N_{L}(r,a;g)$ for $a\in\mathbb{\ol C}$. \end{defi}
\begin{defi}\cite{lahiriNagoya, Lahiri CVEE} Let $f$, $g$ share $(a,0)$. We denote by $\ol N_{*}(r,a;f,g)$ the reduced counting function of those $a$-points of $f$ whose multiplicities differ from the multiplicities of the corresponding $a$-points of $g$. 
\end{defi}
Clearly $$\ol N_{*}(r,a;f,g)\equiv\ol N_{*}(r,a;g,f)\;\;\mbox{and}\;\;\ol N_{*}(r,a;f,g)=\ol N_{L}(r,a;f)+\ol N_{L}(r,a;g).$$ 

\begin{defi} \cite{Lahiri CVEE} For $a\in\mathbb{C}\cup\{\infty\}$ we denote by $N(r,a;f\vline=1)$ the counting function of simple $a$ points of $f$. For a positive integer $m$ we denote by $N(r,a;f\vline\leq m) (N(r,a;f\vline\geq m))$ the counting function of those $a$ points of $f$ whose multiplicities are not greater(less) than $m$ where each $a$ point is counted according to its multiplicity.

$\ol N(r,a;f\vline\leq m)\; (\ol N(r,a;f\vline\geq m))$ are defined similarly, where in counting the $a$-points of $f$ we ignore the multiplicities.

Also $N(r,a;f\vline <m),\; N(r,a;f\vline >m),\; \ol N(r,a;f\vline <m)\; and\; \ol N(r,a;f\vline >m)$ are defined analogously.  \end{defi}
\begin{defi}\cite{ab1} Let $a,b_{1},b_{2},\ldots,b_{q} \in\mathbb{C}\;\cup\{\infty\}$. We denote by $N(r,a;f\vline\; g\neq b_{1},b_{2},\ldots,b_{q})$ the counting function of those $a$-points of $f$, counted according to multiplicity, which are not the $b_{i}$-points of $g$ for $i=1,2,\ldots,q$. \end{defi}
For  two non-constant meromorphic functions $F$ and $G$, 
set\be{\label{e2.2}}H=\left(\frac{F^{''}}{F^{'}}-\frac{2F^{'}}{F-1}\right)-\left(\frac{G^{''}}{G^{'}}-\frac{2G^{'}}{G-1}\right).\ee 


\begin{lem}\label{2.1}\cite{15}
Let $F$, $G$  share $(1,0)$ and $H\not\equiv 0$. Then $$N_{E}^{1)}(r,1;F)=N_{E}^{1)}(r,1;G)\leq N(r,H)+S(r,F)+S(r,G).$$
\end{lem}

\begin{lem}\label{102}
Let $f$ be a non-constant meromorphic function and $\mathcal{L}$ be an non-constant L-function sharing a set $S$ IM, where $|S|\geq3$. Then $\rho(f)=\rho(\mathcal{L})=1$. Furthermore, $ S(r,f)=O(\log r)=S(r,\mathcal{L}).$ 
\end{lem}
\begin{proof}
Proceeding in a similar method as done in the proof of Theorem 5, \cite[see p. 6]{yuan2018}, we would obtain $\rho(f)=\rho(\mathcal{L})=1$. So we omit it.\\
Since $\rho(f)=\rho(\mathcal{L})=1$, so from the definition of $S(r,f)$ we get $S(r,f)=O(\log r)=S(r,\mathcal{L})$.
\end{proof}
\begin{lem}\label{new}
Let $f$, $g$ be two non-constant meromorphic functions sharing (1,t), where $t\in\mathbb{N}\cup\{0\}$. Then
$$\ol N(r,1;f)+\ol N(r,1;g)\leq N_{E}^{1)}(r,1;f)+\left(1-t\right) \ol N_{*}(r,1;f,g)+N(r,1;f).$$

\begin{proof}
Since $f$ and $g$ share $(1,t)$, we observe that $$\ol N(r,1;f)+\ol N(r,1;g)=2\ol N(r,1;f).$$
{\bf Case-I :} Suppose $t\geq 2$.\\
Let $z_0$ be a $1$ point of $f$ with multiplicity $p$ and a $1$ point of $g$ of multiplicity $q$. Since $f$ and $g$ share $(1,t)$, therefore $p\leq t$ implies $p=q$.
\par{\bf Subcase - I :} When $p\leq t$.
If $p=1$, then $z_0$ is counted once in both $N_{E}^{1)}(r,1;f)$ and $N(r,1;f)$. On the other hand $z_0$ is not counted in $\ol N_{*}(r,1;f,g)$. Again if $p\ne 1$, then $z_0$ is counted $p$ times (i.e., at least 2 times) in $N(r,1;f)$ and in this case $z_0$ is not counted in $N_{E}^{1)}(r,1;f) \;\text{and}\; \ol N_{*}(r,1;f,g)$. Therefore $z_0$ is counted at least $2$ times in $N_{E}^{1)}(r,1;f)+\left(1-t\right) \ol N_{*}(r,1;f,g)+N(r,1;f).$
\par{\bf Subcase - II :} When $p\geq (t+1)$. If $p=q$, then $z_0$ is counted $p$ time (i.e., at least 3 times) in  $N(r,1;f)$ and $z_0$ is not counted in $N_{E}^{1)}(r,1;f) \;\text{and}\; \ol N_{*}(r,1;f,g)$. When $p\ne q$, then $z_0$ is counted $(1-t)$ times in $(1-t)\ol N_{*}(r,1;f,g)$ and counted $p$ times (i.e., at least $t+1$ times) in  $N(r,1;f)$ and $z_0$ is not counted in $ N_{E}^{1)}(r,1;f)$; i.e., $z_0$ is counted at least $(1-t)+(t+1)=2$ times in $N_{E}^{1)}(r,1;f)+\left(1-t\right) \ol N_{*}(r,1;f,g)+N(r,1;f).$\\
Now since $z_0$ is counted two times in $\ol N(r,1;f)+\ol N(r,1;g).$ Therefore in any sub case we have 
$$\ol N(r,1;f)+\ol N(r,1;g)\leq N_{E}^{1)}(r,1;f)+\left(1-t\right) \ol N_{*}(r,1;f,g)+N(r,1;f) .$$
{\bf Case-II :} Suppose $t=1$.\\
Then clearly \bea 2\ol N(r,1;f)\nonumber &\leq& N(r,1;f|=1) +N(r,1
;f)\\\nonumber&=&  N_{E}^{1)}(r,1;f)+N(r,1;f).\nonumber \eea
Therefore, for $t=1$
$$\ol N(r,1;f)+\ol N(r,1;g)\leq N_{E}^{1)}(r,1;f)+\left(1-t\right) \ol N_{*}(r,1;f,g)+N(r,1;f).$$

{\bf Case-III :} Suppose that $t=0$.\\
Let $z_0$ be a $1$ point of $f$ with multiplicity $p$ and a $1$ point of $g$ of multiplicity $q$. If $p=q=1$, then $z_0$ is counted $2$ times in both $2\ol N(r,1;f)$ and $N_{E}^{1)}(r,1;f)+\left(1-t\right) \ol N_{*}(r,1;f,g)+N(r,1;f)$, as $\ol N_{*}(r,1;f,g)$ does not count $z_0$. If $p=1,\;q\neq 1$, then also $z_0$ is counted $2$ times in both $2\ol N(r,1;f)$ and\\ $N_{E}^{1)}(r,1;f)+\left(1-t\right)\ol N_{*}(r,1;f,g)+N(r,1;f)$, as in this case $N_{E}^{1)}(r,1;f)$ does not count $z_0$. Finally if $p\neq 1$, then $z_0$ is counted at least 2 times in $\left(1-t\right) \ol N_{*}(r,1;f,g)+N(r,1;f)$ and $z_0$ is not counted in $N_{E}^{1)}(r,1;f)$.\\
Therefore, for $t=0$, $$\ol N(r,1;f)+\ol N(r,1;g)\leq N_{E}^{1)}(r,1;f)+\left(1-t\right) \ol N_{*}(r,1;f,g)+N(r,1;f).$$
And hence in any case,	$$\ol N(r,1;f)+\ol N(r,1;g)\leq N_{E}^{1)}(r,1;f)+\left(1-t\right) \ol N_{*}(r,1;f,g)+N(r,1;f).$$


\end{proof}
\end{lem}

\begin{lem}
Let $F^{*}-1=\d\frac{a_{n}\prod\limits_{i=1}^{n}(f-w_{i})}{\psi(f)}$ and $G^{*}-1=\d\frac{a_{n}\prod\limits_{i=1}^{n}(\mathcal{L}-w_{i})}{\psi(\mathcal
{L})}$, where $f$  be a non-constant meromorphic function, $\mathcal{L}$ be an non-constant L-function, $a_{n},w_{i}\in\mathbb{C}-\{0\}$; $\forall i\in\{1,2,\ldots,n\}$ and  $\psi(z)$ be a polynomial of degree less than $n$ with $\psi(w_{i})\neq0$; $\forall i\in\{1,2,\ldots,n\}$. Further suppose that $F^{*}$ and $G^{*}$  share $(1,t)$, where $t\in\mathbb{N}\cup\{0\}$.  Then 
\be\nonumber \ol{N}_L(r,1;F^{*})\leq \frac{1}{t+1}\left[\ol{ N}(r,\infty;f)+\ol{ N}(r,0;f)-N_1(r,0,f')\right]+O(\log r),\ee
where $N_1(r,0,f')=N(r,0;f'|f\ne 0,w_1,w_2,...,w_n).$
Similar expression also holds for $\ol  N_L(r,1;G^{*})$.
\end{lem}
\begin{proof}
Since $F^{*}$ and $G^{*}$  share $(1,t)$, so using the first fundamental theorem we find that
\bea  \ol{ N}_L(r,1;F^{*}) &\leq& \ol{ N}(r,1;F^{*}|\geq t+2)\nonumber\\&&
\leq \frac{1}{t+1}\left[ N(r,1;F^{*})-\ol{ N}(r,1;F^{*})\right]
\nonumber\\&&
\leq \frac{1}{t+1}\left[\sum_{i=1}^{n}\left(N(r,w_i;f)-\ol{ N}(r,w_i;f)\right)\right]	\nonumber\\&&
\leq \frac{1}{t+1}\left[N(r,0;f'|f\ne 0)-N_1(r,0,f')\right]\nonumber\\&&
\leq \frac{1}{t+1}\left[N(r,0;\frac{f'}{f})-N_1(r,0,f')\right]\nonumber\\&&
\leq \frac{1}{t+1}\left[N(r,\infty;\frac{f}{f^{'}})-N_1(r,0,f')\right]+O(\log r)\nonumber\\&&
\leq \frac{1}{t+1}\left[N(r,\infty;\frac{f^{'}}{f})-N_1(r,0,f')\right]+O(\log r)\nonumber\\&&
\leq \frac{1}{t+1}\left[\ol{ N}(r,\infty;f)+\ol{ N}(r,0;f)-N_1(r,0,f')\right]+O(\log r)\nonumber\
%\\&&
%	\leq \frac{1}{p+1}\left[\ol{ N}(r,0;f)-N_1(r,0,f')\right]+S(r,f).\;\;[\because \ol{ N}(r,\infty;f)=O(\log r)]
\eea %as $f$ has finitely many poles and hence $\ol N(r,\infty;f)=S(r,f)$.
This proves the lemma.
\end{proof}

\begin{lem}\label{l11}

Let $P(z)$, $S$ and $s$ as defined by (\ref{el000}), (\ref{el00}) and (\ref{el0000}) respectively. Suppose that $f$, $\mathcal{L}$ share $(S,t)$, where  $t\in \mathbb{N}\cup\{0\}$ and $f$, $\mathcal{L}$ be a non-constant meromorphic function and an $L$-function respectively. Further suppose that  \be\label{el04}\mathcal{F}=\d\frac{P(f)-a_{0}}{-a_{0}}\;\; and\;\; \mathcal{G}=\d\frac{P(\mathcal{L})-a_{0}}{-a_{0}}.\ee Then for $n\geq 2s+5$, when $t\geq 2$ ,$n\geq 2s+6$, when $t=1$ and for $n\geq2s+11$, when $t=0$ we get the following.
\beas\frac{1}{\mathcal{F}-1}=\frac{A}{\mathcal{G}-1}+B,\eeas where $A(\neq0), B\in\mathbb{C}$.
\end{lem}
\begin{proof}
According to the assumptions of the lemma we clearly have  $\mathcal{F}$, $\mathcal{G}$ share $(1,t)$ and $$\mathcal{F}^{'}=-\frac{na_{n}}{a_{0}}\prod\limits_{i=1}^{s}(f-\eta_{i})^{m_{i}}f^{'};\;\;\; \mathcal{G}^{'}=-\frac{na_{n}}{a_{0}}\prod\limits_{i=1}^{s}(\mathcal{L}-\eta_{i})^{m_{i}}\mathcal{L}^{'},$$ where $\sum\limits_{i=1}^{s}m_{i}=n-1$. 
Now consider $H$ as given by (\ref{e2.2}) for $\mathcal{F}$ and $\mathcal{G}$.
\vspace{0.1in}\par{\bf Case-I:} Suppose $H\not\equiv 0$. Then, it can be easily verified that $H$ has only simple poles and these poles come  from  the following points.
\begin{enumerate}
\item[(i)] $\eta_{i}$ points of $f$ and $\mathcal{L}$.
\item[(ii)] Poles of $f$ and $\mathcal{L}$.
\item[(iii)] $1$ points of $\mathcal{F}$ and $\mathcal{G}$ having different multiplicities.
\item[(iv)] Those zeros of $f^{'}$ and $\mathcal{L}^{'}$ which are not zeros of   $\prod\limits_{i=1}^{s}(f-\eta_{i})(\mathcal{F}-1)$ and $\prod\limits_{i=1}^{s}(\mathcal{L}-\eta_{i})(\mathcal{G}-1)$ respectively.
\end{enumerate}
Therefore we obtain 
\bea\label{e2.1} N(r,H)&\leq& \sum\limits_{i=1}^{s}\left[\ol N(r,\eta_{i};f)+\ol N(r,\eta_{i};\mathcal{L})\right]+\ol N(r,\infty;f)+\ol N(r,\infty;\mathcal{L})\\\nonumber&&+\ol N_{*}(r,1;\mathcal{F},\mathcal{G})+\ol N_{0}(r,0,f^{'})+\ol N_{0}(r,0,\mathcal{L}^{'}),\eea where $\ol N_{0}(r,0,f^{'})$ and $\ol N_{0}(r,0,\mathcal{L}^{'})$ denotes the reduced counting functions of those zeros of $f^{'}$ and $\mathcal{L^{'}}$ which are not zeros of   $\prod\limits_{i=1}^{s}(f-\eta_{i})(\mathcal{F}-1)$ and $\prod\limits_{i=1}^{s}(\mathcal{L}-\eta_{i})(\mathcal{G}-1)$ respectively.
Using the second fundamental theorem we get 
\bea\label{e2.3}&&(n+s-1)T(r,f)\\\nonumber&\leq&\ol N(r,1;\mathcal{F})+\sum\limits_{i=1}^{s}\ol N(r,\eta_{i};f)+\ol N(r,\infty;f)-N_{0}(r,0,f^{'})+S(r,f),\eea
\bea\label{e2.4}&&(n+s-1)T(r,\mathcal{L})\\\nonumber&\leq&\ol N(r,1;\mathcal{G})+\sum\limits_{i=1}^{s}\ol N(r,\eta_{i};\mathcal{L})+\ol N(r,\infty;\mathcal{L})-N_{0}(r,0,\mathcal{L}^{'})+S(r,\mathcal{L}).\eea
Let us denote by $T(r)=\text{max}\{T(r,f),T(r,\mathcal{L})\}$. Now combining (\ref{e2.3}) and (\ref{e2.4}) and using {\em Lemma} {\bf\ref{102}} and the fact that $\ol{N}(r,\infty;\mathcal{L})=O(\log r)$, we get
\bea
\label{e2.5}&&(n+s-1)\{T(r,f)+T(r,\mathcal{L})\}\\\nonumber&\leq&\ol N(r,1;\mathcal{F})+\ol N(r,1;\mathcal{G})+\sum\limits_{i=1}^{s}\left[\ol N(r,\eta_{i};f)+\ol N(r,\eta_{i};\mathcal{L})\right]\\\nonumber&&+\ol N(r,\infty;f)-N_{0}(r,0,f^{'})-N_{0}(r,0,\mathcal{L}^{'})+O(\log r).\nonumber
\eea

%From \ref{e2.2} we obtain \be\label{eq111} m(r,H)=S(r) \ee.

Now it follows from {\em Lemma} {\bf\ref{new}}, {\em Lemma} {\bf\ref{102}}, {\em Lemma} {\bf\ref{2.1}}, (\ref{e2.1}) and in view of $\ol{N}(r,\infty;\mathcal{L})=O(\log r)$ that 
\bea\label{eq11.4}&& \ol N(r,1;\mathcal{F})+\ol N(r,1;\mathcal{G})\\\nonumber&\leq&  (s+1)T(r,f)+sT(r,\mathcal{L})+(2-t)\ol N_{*}(r,1;\mathcal{F},\mathcal{G})+N(r,1;\mathcal{F})\\\nonumber&&+\ol N_{0}(r,0,f^{'})+\ol N_{0}(r,0,\mathcal{L}^{'})+O(\log r).\eea
and \bea\label{eq11.5}&&\ol N(r,1;\mathcal{F})+\ol N(r,1;\mathcal{G})\\\nonumber&\leq&  (s+1)T(r,f)+sT(r,\mathcal{L})+(2-t)\ol N_{*}(r,1;\mathcal{F},\mathcal{G})+N(r,1;\mathcal{G})\\\nonumber&&+\ol N_{0}(r,0,f^{'})+\ol N_{0}(r,0,\mathcal{L}^{'})+O(\log r).\eea\\

Now from (\ref{e2.5}) , (\ref{eq11.4}) and using $T(r,\mathcal{F})=nT(r,f)$ we get
\bea(n+s-1)\{T(r,f)+T(r,\mathcal{L})\}\nonumber&\leq&(2s+2)T(r,f)+2sT(r,\mathcal{L})\\\nonumber&&+(2-t)\ol N_{*}(r,1;\mathcal{F},\mathcal{G})+nT(r,f)+O(\log r);\eea
i.e.,
\bea\label{eq11.6}&& nT(r,\mathcal{L})\\\nonumber
&\leq& (s+3)T(r,f)+(s+1)T(r,\mathcal{L})+(2-t)\ol N_{*}(r,1;\mathcal{F},\mathcal{G})+O(\log r).\\\nonumber&
%\leq& (s+3)T(r,f)+(s+1)T(r,\mathcal{L})+\\\nonumber&& \frac{(2-t)}{t+1}\left[\ol{ N}(r,\infty;f)+\ol{ N}(r,0;f)+\ol{ N}(r,\infty;\mathcal{L})+\ol{ N}(r,0;\mathcal{L})\right]+S(r) \nonumber\\&\leq& (s+3+\frac{2(2-t)}{t+1})T(r,f)+(s+1+\frac{(2-t)}{t+1})T(r,\mathcal{L})+S(r)\nonumber\\&\leq& \left(2s+4+\frac{3(2-t)}{t+1}\right)T(r)+S(r)
.\nonumber\eea

Also from (\ref{e2.5}) , (\ref{eq11.5}) and using $T(r,\mathcal{G})=nT(r,\mathcal{L})$ we get

\bea(n+s-1)\{T(r,f)+T(r,\mathcal{L})\}\nonumber&\leq&(2s+2)T(r,f)+2sT(r,\mathcal{L})+\\\nonumber&&(2-t)\ol N_{*}(r,1;\mathcal{F},\mathcal{G})+nT(r,\mathcal{L})+O(\log r);\eea
i.e.,
\bea\label{eq11.7}&& nT(r,f)\\\nonumber
&\leq& (s+3)T(r,f)+(s+1)T(r,\mathcal{L})+(2-t)\ol N_{*}(r,1;\mathcal{F},\mathcal{G})+O(\log r).\\\nonumber&
%\leq& (s+3)T(r,f)+(s+1)T(r,\mathcal{L})+\\\nonumber&& \frac{(2-t)}{t+1}\left[\ol{ N}(r,\infty;f)+\ol{ N}(r,0;f)+\ol{ N}(r,\infty;\mathcal{L})+\ol{ N}(r,0;\mathcal{L})\right]+S(r) \nonumber\\&\leq& (s+3+\frac{2(2-t)}{t+1})T(r,f)+(s+1+\frac{(2-t)}{t+1})T(r,\mathcal{L})+S(r)\nonumber\\&\leq& \left(2s+4+\frac{3(2-t)}{t+1}\right)T(r)+S(r)
.\nonumber\eea
Therefore, from (\ref{eq11.6}) and (\ref{eq11.7}) we get,

\be\label{11.8}nT(r)\leq (2s+4)T(r)+(2-t)\ol N_{*}(r,1;\mathcal{F},\mathcal{G})+O(\log r).\ee
Now for $t\geq2$; from (\ref{11.8}) we get$$nT(r)\leq (2s+4)T(r)+O(\log r),$$
which is a contradiction for $n\geq 2s+5$.\\          
When $t <2$, from (\ref{11.8}) we get
\bea\label{eq11.9} nT(r)
&\leq& (2s+4)T(r)+(2-t)\ol N_{*}(r,1;\mathcal{F},\mathcal{G})+O(\log r)\\\nonumber&
\leq& (2s+4)T(r)+\frac{(2-t)}{t+1}[\ol{ N}(r,\infty;f)+\ol{ N}(r,0;f)\\\nonumber&&+\ol{ N}(r,\infty;\mathcal{L})+\ol{ N}(r,0;\mathcal{L})]+O(\log r) \nonumber\\&\leq& \left(2s+4+\frac{3(2-t)}{t+1}\right)T(r)+O(\log r)
.\nonumber\eea

Now for $t=1$; from (\ref{eq11.9}) we get$$nT(r)\leq (2s+4+\frac{3}{2})T(r)+O(\log r),$$
which is a contradiction for $n\geq 2s+6$.\\  
For $t=0$; from (\ref{eq11.9}) we get $$nT(r)\leq (2s+4+6)T(r)+O(\log r),$$
which is a contradiction for $n\geq 2s+11$.\\



\vspace{0.1in}\par{\bf Case-II:} Suppose	$H\equiv0$. Now integrating (\ref{e2.2}), we find that \bea\nonumber \frac{1}{\mathcal{F}-1}=\frac{A}{\mathcal{G}-1}+B,\;\;\; \mbox{where}\;\; A(\neq0), B\in\mathbb{C}.\eea 	
\end{proof}

\begin{lem}\label{l1.1}
Let $R(z)$, $S$ and $l$ as defined by (\ref{el2}), (\ref{el02}) and (\ref{el3}) respectively. Suppose that $f$, $\mathcal{L}$ share $(S,t)$, where  $t\in \mathbb{N}\cup\{0\}$ and $f$, $\mathcal{L}$ be a non-constant meromorphic function and an $L$-function respectively. Further suppose that \be\label{el05}\mathbb{F}= R(f)\;\; and\;\; \mathbb{G}=R(\mathcal{L}).\ee 
Then for $n\geq 2l+5$, when $t\geq2$, for $n\geq 2l+6$, when $t=1$ and for $n\geq 2l+11$, when $t=0$ we get the following.
\beas\frac{1}{\mathbb{F}-1}=\frac{A}{\mathbb{G}-1}+B,\eeas where $A(\neq0), B\in\mathbb{C}$.
\end{lem}
\begin{proof} 	
Clearly $\mathbb{F}, \mathbb{G}$ share $(1,t)$ and in view of (\ref{el3}) we have
\be \nonumber \mathbb{F}^{'}=\frac{\gamma\prod\limits_{j=1}^{l}(f-\delta_j)^{q_{j}}}{\prod\limits_{j=1}^{k}(f-\beta_j)^{p_{j}}}f',\;\;\;\;\;\;\mathbb{G}^{'}=\frac{\gamma\prod\limits_{j=1}^{l}(\mathcal{L}-\delta_j)^{q_{j}}}{\prod\limits_{j=1}^{k}(\mathcal{L}-\beta_j)^{p_{j}}}\mathcal{L}'.\ee %where $\gamma\in\mathbb{C}-\{0\}$, $\delta_{j}$ be the distinct zeros of $R^{'}(z)$ and $q_{j}$ being their respective multiplicities for all $j\in\{1,2,\ldots,l\}$.
Now consider $H$ as given by (\ref{e2.2}) for $\mathbb{F}$ and $\mathbb{G}$.
\vspace{0.1in}\par{\bf Case-I:} Suppose $H\not\equiv 0$. Since $H$ has only simple poles and in this case these poles come  from  the following points.
\begin{enumerate}
\item[(i)] $\delta_{j}$ -points of $f$ and $\mathcal{L}$.
\item[(ii)] Poles of $f$ and $\mathcal{L}$.
\item[(iii)] $1$ points of $\mathbb{F}$ and $\mathbb{G}$ having different multiplicities.
\item[(iv)] Those zeros of $f^{'}$ and $\mathcal{L^{'}}$ which are not zeros of   $\prod\limits_{j=1}^{l}(f-\delta_{j})(\mathbb{F}-1)$ and $\prod\limits_{j=1}^{l}(\mathcal{L}-\delta_{j})(\mathbb{G}-1)$ respectively.
\end{enumerate}
Therefore we obtain 
\bea\label{ep2} N(r,H) &\leq& \ol{N}(r,\infty;f)+\sum_{j=1}^{l}\ol{N}(r,\delta_j;f)+\ol N_0(r,0;f') +\ol{N}(r,\infty;\mathcal{L})\\\nonumber&&+\sum_{j=1}^{l}\ol{N}(r,\delta_j;\mathcal{L})+\ol N_0(r,0;\mathcal{L}')+\ol N_*(r,1;\mathbb{F},\mathbb{G})\\\nonumber&&+S(r,f)+S(r,\mathcal{L}),\eea
where we write $\ol N_0(r,0;f')$ for the reduced counting function of the zeros of $f'$ that are not zeros
of $(\mathbb{F}-1)\prod\limits_{j=1}^{l}(f-\delta_j)^{q_{j}}$ and $\ol N_0(r,0;\mathcal{L}')$ is similarly defined. Now using the second fundamental theorem we get
\bea\label{e3.3}&&(n+l-1)T(r,f)\\\nonumber&\leq&\ol N(r,1;\mathbb{F})+\sum\limits_{i=1}^{l}\ol N(r,\delta_{i};f)+\ol N(r,\infty;f)-N_{0}(r,0,f^{'})+S(r,f),\eea
\bea\label{e3.4}&&(n+l-1)T(r,\mathcal{L})\\\nonumber&\leq&\ol N(r,1;\mathbb{G})+\sum\limits_{i=1}^{l}\ol N(r,\delta_{i};\mathcal{L})+\ol N(r,\infty;\mathcal{L})-N_{0}(r,0,\mathcal{L}^{'})+S(r,\mathcal{L}).\eea
Let us denote by $T(r)=\text{max}\{T(r,f),T(r,\mathcal{L})\}$. 	 Now combining (\ref{e3.3}), (\ref{e3.4}) and using {\em Lemma} {\bf\ref{102}} and the fact that $\ol{N}(r,\infty;\mathcal{L})=O(\log r)$, we get
\bea\label{e3.5}&&(n+l-1)\{T(r,f)+T(r,\mathcal{L})\}\\\nonumber&\leq&\ol N(r,1;\mathbb{F})+\ol N(r,1;\mathbb{G})+\sum\limits_{i=1}^{l}\left[\ol N(r,\delta_{i};f)+\ol N(r,\delta_{i};\mathcal{L})\right]\\\nonumber&&+\ol N(r,\infty;f)-N_{0}(r,0,f^{'})-N_{0}(r,0,\mathcal{L}^{'})+O(\log r).\eea 	  		
%From \ref{e2.2} we obtain \be\label{eq111} m(r,H)=S(r) \ee.
Now it follows from {\em Lemma} {\bf\ref{new}}, {\em Lemma} {\bf\ref{102}}, {\em Lemma} {\bf\ref{2.1}} and (\ref{ep2}) and in view of $\ol{N}(r,\infty;\mathcal{L})=O(\log r)$ that 
\bea \label{eq12.4} &&\ol N(r,1;\mathbb{F})+\ol N(r,1;\mathbb{G})\\\nonumber&\leq&  (l+1)T(r,f)+lT(r,\mathcal{L})+(2-t)\ol N_{*}(r,1;\mathbb{F},\mathbb{G})+N(r,1;\mathbb{F})\\\nonumber&&+\ol N_{0}(r,0,f^{'})+\ol N_{0}(r,0,\mathcal{L}^{'})+O(\log r).\eea
and \bea\label{eq12.5}&& \ol N(r,1;\mathbb{F})+\ol N(r,1;\mathbb{G})\\\nonumber&\leq& (l+1)T(r,f)+lT(r,\mathcal{L})+(2-t)\ol N_{*}(r,1;\mathbb{F},\mathbb{G})\\\nonumber&&+N(r,1;\mathbb{G})+\ol N_{0}(r,0,f^{'})+\ol N_{0}(r,0,\mathcal{L}^{'})+O(\log r).\eea 	  	
Now from (\ref{e3.5}), (\ref{eq12.4}) and using $T(r,\mathbb{F})=nT(r,f)$ we get
\bea(n+l-1)\{T(r,f)+T(r,\mathcal{L})\}\nonumber&\leq&(2l+2)T(r,f)+2lT(r,\mathcal{L})+(2-t)\ol N_{*}(r,1;\mathbb{F},\mathbb{G})\\\nonumber&&+nT(r,f)+O(\log r);\eea
i.e.,
\bea\label{eq12.6} &&nT(r,\mathcal{L})\\\nonumber
&\leq& (l+3)T(r,f)+(l+1)T(r,\mathcal{L})+(2-t)\ol N_{*}(r,1;\mathbb{F},\mathbb{G})+O(\log r).\\\nonumber&
%\leq& (l+3)T(r,f)+(l+1)T(r,\mathcal{L})+\\\nonumber&& \frac{(2-t)}{t+1}\left[\ol{ N}(r,\infty;f)+\ol{ N}(r,0;f)+\ol{ N}(r,\infty;\mathcal{L})+\ol{ N}(r,0;\mathcal{L})\right]+S(r) \nonumber\\&\leq& (l+3+\frac{2(2-t)}{t+1})T(r,f)+(l+1+\frac{(2-t)}{t+1})T(r,\mathcal{L})+S(r)\nonumber\\&\leq& \left(2l+4+\frac{3(2-t)}{t+1}\right)T(r)+S(r).
\nonumber\eea 	   
Also from (\ref{e3.5}), (\ref{eq12.5}) and using $T(r,\mathbb{G})=nT(r,\mathcal{L})$ we get 	   
\bea(n+l-1)\{T(r,f)+T(r,\mathcal{L})\}\nonumber&\leq&(2l+2)T(r,f)+2lT(r,\mathcal{L})+(2-t)\ol N_{*}(r,1;\mathbb{F},\mathbb{G})\\\nonumber&&+nT(r,\mathcal{L})+O(\log r);\eea
i.e.,
\bea\label{eq12.7}&& nT(r,f)\\\nonumber
&\leq& (l+3)T(r,f)+(l+1)T(r,\mathcal{L})+(2-t)\ol N_{*}(r,1;\mathbb{F},\mathbb{G})+O(\log r).\\\nonumber&
%\leq& (l+3)T(r,f)+(l+1)T(r,\mathcal{L})+\\\nonumber&& \frac{(2-t)}{t+1}\left[\ol{ N}(r,\infty;f)+\ol{ N}(r,0;f)+\ol{ N}(r,\infty;\mathcal{L})+\ol{ N}(r,0;\mathcal{L})\right]+S(r) \nonumber\\&\leq& (l+3+\frac{2(2-t)}{t+1})T(r,f)+(l+1+\frac{(2-t)}{t+1})T(r,\mathcal{L})+S(r)\nonumber\\&\leq& \left(2l+4+\frac{3(2-t)}{t+1}\right)T(r)+S(r).
\nonumber\eea
Therefore, from (\ref{eq12.6}) and (\ref{eq12.7}) we get,

\be\label{12.8}nT(r)\leq (2l+4)T(r)+(2-t)\ol N_{*}(r,1;\mathbb{F},\mathbb{G})+O(\log r).\ee
Now for $t\geq2$; from (\ref{12.8}) we get$$nT(r)\leq (2l+4)T(r)+O(\log r),$$
which is a contradiction for $n\geq 2l+5$.\\          
When $t <2$, from (\ref{12.8}) we get
\bea\label{eq12.9} &&nT(r)\\\nonumber
&\leq& (2l+4)T(r)+(2-t)\ol N_{*}(r,1;\mathbb{F},\mathbb{G})+O(\log r)\\\nonumber&
\leq& (2l+4)T(r)+ \frac{(2-t)}{t+1}\ol{ N}(r,\infty;f)+\ol{ N}(r,0;f)+\ol{ N}(r,\infty;\mathcal{L})\\\nonumber&&+\ol{ N}(r,0;\mathcal{L})+O(\log r) \nonumber\\&\leq& \left(2l+4+\frac{3(2-t)}{t+1}\right)T(r)+O(\log r)
.\nonumber\eea

Now for $t=1$; from (\ref{eq12.9}) we get$$nT(r)\leq (2l+4+\frac{3}{2})T(r)+O(\log r),$$
which is a contradiction for $n\geq 2l+6$.\\  
For $t=0$; from (\ref{eq12.9}) we get $$nT(r)\leq (2l+4+6)T(r)+O(\log r),$$
which is a contradiction for $n\geq 2l+11$.\\
% Therefore, from (\ref{eq12.6}) and (\ref{eq12.7}) we get,

% \be\label{12.8}nT(r)\leq \left(2l+4+\frac{3(2-t)}{t+1}\right)T(r)+S(r).\ee
%  Now for $t=0$; from (\ref{12.8}) we get $$nT(r)\leq (2l+4+6)T(r)+S(r),$$
%       which is a contradiction for $n\geq 2l+11$.\\
%  For $t=1$; from (\ref{12.8}) we get$$nT(r)\leq (2l+4+\frac{3}{2})T(r)+S(r),$$
%      which is a contradiction for $n\geq 2l+6$.\\
%  For $t\geq2$; from (\ref{12.8}) we get$$nT(r)\leq (2l+4)T(r)+S(r),$$
%      which is a contradiction for $n\geq 2l+5$.

\vspace{0.1in}\par{\bf Case-II:} Suppose	$H\equiv0$. Now integrating (\ref{e2.2}), we find that \bea\nonumber \frac{1}{\mathbb{F}-1}=\frac{A}{\mathbb{G}-1}+B,\;\;\; \mbox{where}\;\; A(\neq0), B\in\mathbb{C}.\eea 
%	i.e; \beas\frac{\phi(f)}{P(f)}=\frac{A\phi(\mathcal{L})}{P(\mathcal{L})}+B,\;\;\; \mbox{where}\;\; A(\ne0),B(=-A^{'})\in\mathbb{C}.\eeas
\end{proof}

\begin{lem}\label{l2}\cite{hxy0}
Let $F$ and $G$ be two non-constant meromorphic functions such that\\ \beas\displaystyle\frac{1}{F-1}=\frac{A}{G-1}+B,\eeas where $A(\ne0),B\in\mathbb{C}$. If
\begin{equation*}
\ol{N}(r,0; F)+	\ol{N}(r,\infty; F)+	\ol{N}(r,0; G)+	\ol{N}(r,\infty; G) < T(r),
\end{equation*} 
where $T(r)=\max\{T(r,F),T(r,G)\}$. Then either $FG=1$ or $F=G.$
\end{lem}
\begin{lem}\label{14}
%	Let $F=\prod\limits_{i=1}^{p}(f-\alpha_i)^{m_i}$ and $G=\prod\limits_{i=1}^{p}(\mathcal{L}-\alpha_i)^{m_i}$, where $f$ be a non-constant meromorphic function having finitely number of poles and $\mathcal{L}$ be an L-function, $p \geq 2$, $\sum\limits_{i=1}^{p}m_i=n$ and $\alpha_i \in \mathbb{C}$ for $i=1, 2, \ldots ,p$. 
Let $\mathcal{F}$, $\mathcal{G}$ be defined by (\ref{el04}). If $p \geq 2$
%\begin{enumerate}
%  \item[(i)] $p\geq 3$, or
% \item [(ii)] $p=2$ and $gcd(m_i,n)=1$ for at least one of the $m_i$'s such that $n=\sum\limits_{i=1}^{2} m_i \geq 3$, or
%\item[(c)] $p=2$ and $gcd(m_i,n)\ne 1$ for each $m_i$ such that $n\geq (b_1+b_2)+1$, where $b_i=gcd(m_i,n)$,
%\end{enumerate}
Then $\mathcal{F}\mathcal{G} \not= a$, where $a$ is non-zero complex constant.  
\end{lem}
\begin{proof}
%	For the proof of this lemma we proceed in a similar way as the proof of {\em Lemma \ref{l3}} except part (ii) and part (iii) which are clarified below. 
%	\par Let $FG=a$.
%	\par For $k\geq 3$, the proof is same as the proof of part (i) of {\em Lemma \ref{l3}}.
On the contrary, suppose that $\mathcal{F}\mathcal{G} = a$. Then
\be\label{e01}\prod_{i=1}^{p}(f-\alpha_i)^{m_i}\prod_{i=1}^{p}(\mathcal{L}-\alpha_i)^{m_i}=a\left(\frac{a_{0}}{a_{n}}\right)^{2}=a_{1}(say).\ee

It is clear from (\ref{e01} ) that each $\alpha_i$-point of $f$ is a pole of $\mathcal{L}$ and vice-versa.\\
%Now let us consider the following cases.%We prove this part in two sub-parts as follows.\\ 
{\bf\underline{Case-i:}} Let $p\geq 4$. Since  an L- function has at most one pole, then in view of (\ref{e01}) we can say that $f$ has at least three $\alpha_i$-points which are picard exeptional values. That is, the meromorphic function $f$ omits at least 3 values, so  $f$ must be constant. This contradicts our assumption. \\
{\bf\underline{Case-ii:}} Let $p = 3$. Again like the arguments made above we can say that $f$ omits two values say $\alpha_{1}, \alpha_{2}$. Hence $\alpha_{3}$ points of $f$ are the poles of $\mathcal{L}$. Since $z=1$ is the only pole of $\mathcal{L}$, so let  $z=1$ be the $\alpha_3$-point of $f$ of multiplicity $r$ and the pole of $\mathcal{L}$ of multiplicity $u$. Then $m_3r = nu$. Since $u\geq 1$, so $r \geq \frac{n}{m_{3}}$; i.e., $\frac{1}{r}\leq \frac{m_{3}}{n}$.
Hence using the second fundamental theorem, we obtain 
\begin{flalign*}
T(r, f) & \leq \sum_{i=1}^{3}\ol{N}(r, \alpha_i; f)+S(r,f)\\
& \leq \frac{m_3}{n}{N}(r,\alpha_3; f)+O(\log r),
\end{flalign*}
which is  a  contradiction as $m_3 < n$. \\
{\bf\underline{Case-iii}:} Let $p=2$.\\ Since  an L- function has at most one pole, then in view of (\ref{e01}) we can say that $f$ omits at least one $\alpha_i$-points, say $\alpha_1$.

%and $gcd(m_1,n)=1$. Since $m_1+m_2=n$, so $gcd(m_2,n)=1$. % we have the following analysis.

Let $z_0$ be a $\alpha_2$ point of $\mathcal{L}$ of multiplicity $r$, then it will be a pole of $f$ of multiplicity $\nu$, such that $rm_2=\nu n$. Since $\nu \geq 1$, so $r \geq \frac{n}{m_2}$ ; i.e., $\frac{1}{r} \leq\frac{m_2}{n}.$ Now using the second fundamental theorem  in view of $\ol{N}(r,\infty;\mathcal{L})=O(\log r)$ we get
\beas T(r,\mathcal{L}) &\leq& \ol N(r,\alpha_{1};\mathcal{L})+\ol N(r,\alpha_{2};\mathcal{L})+\ol  N(r,\infty;\mathcal{L})+S(r,\mathcal{L})\\&\leq& \frac{m_2}{n} T(r,\mathcal{L})+O(\log r),\eeas
which is a contradiction as $n>m_2$. 	
%(iii) Let $p=2$ and $gcd(m_i,n)=b_i\ne 1$ for $i=1,2$. Let $z_0$ be a $\alpha_i$ point of $\mathcal{L}$ of multiplicity $\xi_i$ and a pole a $f$ of multiplicity $\nu_i$. Then $m_i \xi_i =\nu_i n$. Since $gcd(m_i,n)=b_i$, so $m_i=b_i \eta_i$ and $n=b_i \zeta_i$ for some $\eta_i,\zeta_i \in \mathbb{N}$, where $gcd(\eta_i,\zeta_i)=1$. Now $m_i \xi_i =\nu_i n$ implies $b_i\eta_i \xi_i=\nu_i b_i \zeta_i$; i.e.,  $\eta_i \xi_i=\nu_i \zeta_i$. Since $gcd(\eta_i,\zeta_i)=1$, so $\zeta_i$ divides $\xi_i$. Therefore $\xi_i \geq \zeta_i$; i.e., $\frac{1}{\xi_i}\leq \frac{1}{\zeta_i} = \frac{b_i}{n}$. Now using the second fundamental theorem we get 
%\beas T(r,\mathcal{L}) &\leq& \ol N(r,\alpha_{1};\mathcal{L})+\ol N(r,\alpha_{2};\mathcal{L})+\ol  N(r,\infty;\mathcal{L})+S(r,\mathcal{L})\\&\leq& \frac{(b_1+b_2)}{n} T(r,f)+S(r,\mathcal{L}),\eeas
%	which is a contradiction as $n\geq (b_1+b_2)+1$.
\end{proof}

\begin{lem}\label{l04}
%	Let $F=\d\frac{f^n}{\prod\limits_{j=1}^{k}(f-\beta_j)^{m_{j}}}$ and $G=\d\frac{\mathcal{L}^n}{\prod\limits_{j=1}^{k}(\mathcal{L}-\beta_j)^{m_{j}}}$, where $f$ be a non-constant meromorphic function having finite number of poles and $\mathcal{L}$ be an L-function, $n\geq 2$, $\sum\limits_{j=1}^{k}m_{j}\leq n-1$ and $\beta_j\in \mathbb{C}-\{0\}$ for $j=1,2,\ldots ,k$. 
Let $\mathbb{F}, \mathbb{G}$ as defined by (\ref{el05}). If
\begin{enumerate}
\item[(a)] $k\geq2$, or
\item[(b)] $k=1$ and $n\geq 3b_1+1$, where $b_1=gcd(m_1,n)$,
\end{enumerate} 
Then $\mathbb{F}\mathbb{G} \not= a$, where $a$ is non-zero complex constant.\end{lem}
\begin{proof}
On the contrary suppose that $\mathbb{F} \mathbb{G}\equiv a$. Then 
\bea\label{000} \d\frac{f^n}{\prod\limits_{j=1}^{k}(f-\beta_j)^{m_{j}}}.\d\frac{\mathcal{L}^n}{\prod\limits_{j=1}^{k}(\mathcal{L}-\beta_j)^{m_{j}}}\equiv a\left(\frac{a_{i}}{a_{n}}\right)^{2}=a^{'}(say),\eea and hence   \bea\label{0000} T(r,f)= T(r,\mathcal{L})+O(1).\eea 		
It is clear from (\ref{000} ) that  $\beta_j$ point of $f$ is a zero  of $\mathcal{L}$ and vice-versa. 

{\bf(a)} Let $k\geq 2$.
If $z_0$ be a zero of $\mathcal{L}-\beta_j$ with
multiplicity $p$, then $z_0$ is a zero of $f$ with multiplicity $q$ such that $m_jp=nq$; i.e., $p\d\geq \frac{n}{m_j}$. Therefore  $\ol{N}(r,\beta_j;\mathcal{L})\leq \d\frac{m_j}{n}N(r,\beta_j;\mathcal{L})$. %Again $f$ has only finite number of poles so  $\ol{N}(r,\infty;f)=S(r,f)$

So, using the second fundamental theorem and $\ol{N}(r,\mathcal{L})=O(\log r)$, we get \bea\nonumber (k-1)T(r,\mathcal{L})&\leq& \sum_{j=1}^{k}\ol{N}(r,\beta_j;\mathcal{L})+\ol{N}(r,\infty;\mathcal{L})+S(r,\mathcal{L}) \nonumber\\ & 
\leq& \sum_{j=1}^{k}\frac{m_j}{n}T(r,\mathcal{L})+O(\log r)\nonumber\\&
\leq& (1-\frac{1}{n})T(r,\mathcal{L})+O(\log r)\nonumber,
\eea
which  contradicts $k\geq2$.

{\bf(b)} Let $k=1$ and $gcd(n,m_1)=b_1$, therefore $gcd(n,n-m_1)=b_1$. Then from (\ref{000}) we have 
\bea\label{ec2} \frac{\mathcal{L}^n}{(\mathcal{L}-\beta_1)^{m_1}}=\frac{a^{'}(f-\beta_1)^{m_1}}{f^n}. \eea\\
If $z_0$ be a zero of $\mathcal{L}-\beta_1$ with
multiplicity $p$, then $z_0$ is a zero of $f$ with multiplicity $q$ such that $m_1p=nq$. Since $gcd(n,m_1)=b_1$, so $m_1=b_1 \eta$ and $n=b_1 \zeta$ for some $\eta,\zeta \in \mathbb{N}$, where $gcd(\eta,\zeta)=1$. Now $m_1 p = nq$ implies $b_1\eta p=q b_1\zeta$; i.e.,  $\eta p=q\zeta$. Since $gcd(\eta,\zeta)=1$, so $\zeta$ divides $p$. Therefore $p \geq \zeta$; i.e., $p\d\geq \frac{n}{b_1}$. Therefore, $\ol{N}(r,\beta_1;\mathcal{L})\leq \frac{b_1}{n}T(r,\mathcal{L})$ and similarly $\ol{N}(r,\beta_1;f)\leq \frac{b_1}{n}T(r,f)$ .\\
Also let $z_0$ be a pole of $f$ of multiplicity $\nu$, then $z_0$ will be a zero of $\mathcal{L}$ of multiplicity $\xi$ such that $\nu (n-m_1) = \xi n$. Since $gcd(n,n-m_1)=b_1$, hence $\frac{1}{\nu}\leq \frac{b_1}{n} $. Therefore $\ol{N}(r,\infty;f)\leq \frac{b_1}{n}T(r,f).$ Also we know that $\ol{N}(r,\infty;\mathcal{L})=O(\log r).$\\
Now from (\ref{ec2}) it is clear that $\ol{N}(r,0;\mathcal{
L})= \ol{N}(r,\infty;f)+\ol{N}(r,\beta_1;f).$ \\
Now using the second fundamental theorem in view of (\ref{0000}) we get \bea T(r,\mathcal{L})&\leq& \ol{N}(r,\beta_1;\mathcal{L})+\ol{N}(r,\infty;\mathcal{L})+\ol{N}(r,0;\mathcal{L})+S(r,\mathcal{L}) \nonumber\\ & 
\leq& \frac{3b_1}{n}T(r,\mathcal{L})+O(\log r)\nonumber
\eea
Which is a contradiction as $n\geq 3b_1+1$.
\end{proof}

\section{Proof Of The Theorems}

{\bf Proof of  Theorem 1.1 :} 
We  prove the theorem step by step as follows.

\underline{$\bf{(i) \implies (ii)}$ :
} Suppose $f$ be a non-constant meromorphic function and $\mathcal{L}$ be a non-constant L-function such that  %$P(f)=P(\mathcal{L})\implies f=\mathcal{L}$ and
$E_f(S,t)=E_{\mathcal{L}}(S,t)$, where $t\in\mathbb{N}\cup\{0\}$. Consider $\mathcal
{F}$ and $\mathcal
{G}$ as defined by (\ref{el04}). Then for \begin{itemize}
\item [(a)]$t\geq 2$ and $n\geq 2s+5$, or
\item [(b)] $t=1$ and $n\geq 2s+6$, or
\item [(c)]$t=0$ and $n\geq 2s+11$,
\end{itemize} in view of the {\em Lemma \ref{l11}} we get 
$\d\frac{1}{\mathcal
{F}-1}=\frac{A}{\mathcal
{G}-1}+B$, where $A(\neq0), B\in \mathbb{C}$.
Hence we have \be\label{e1}T(r,\mathcal{F})=T(r,\mathcal{G})+O(1).\ee Since  \be\label{e3} T(r,\mathcal{F})=nT(r,f)+O(1)\;\; and\;\; T(r,\mathcal{G})=nT(r,\mathcal{L})+O(1).\ee So (\ref{e1}) implies that
\be\label{e2}T(r,f)=T(r,\mathcal{\mathcal{L}})+O(1).\ee
Now in view of $\ol{N}(r,\infty,\mathcal{L})=O(\log r)$ using (\ref{e3}) and (\ref{e2}) we get 
\beas &&
\ol{N}(r,0; \mathcal{F})+	\ol{N}(r,\infty; \mathcal{F})+	\ol{N}(r,0; \mathcal{G})+	\ol{N}(r,\infty; \mathcal{G})\\ &\leq& pT(r,f)+pT(r,\mathcal{L})+\ol N(r,\infty;f)+\ol N(r,\infty;\mathcal{L})\\
&=&(2p+1)T(r,f)+O(\log r)\\
&<&\frac{2p+2}{n}T(r,\mathcal{F})\\&\leq& T(r,\mathcal{F}) \;\;\; [\because n \geq 2p + 2].
\eeas
So in view of {\em Lemma \ref{l2}}, we have either $\mathcal{F}\mathcal{G} = 1$ or $\mathcal{F} = \mathcal{G}$.

Now again in view of {\em Lemma \ref{14}} we have $\mathcal{F} \mathcal{G}\neq 1$. Hence $\mathcal{F}= \mathcal{G}$. That is, we get 
\be\nonumber
P(f)=
P(\mathcal{L}),
\ee
which by condition (i) implies $f=\mathcal{L}$.\\

\underline{$\bf{(ii)\implies (i)}$ :}
Let $P(f)=P(\mathcal{L})$. That is, 
$$\prod_{i=1}^{p}(f-\alpha_i)^{m_i}=\prod_{i=1}^{p}(\mathcal{L}-\alpha_i)^{m_i},$$ which implies $f$ and $\mathcal{L}$ share $(S,\infty)$. Therefore, obviously $f$ and $\mathcal{L}$ share $(S,t)$ for $t\in\mathbb{N}\cup\{0\}$. Hence by condition (ii), we have $f= \mathcal{L}$. \\

%\underline{$\bf{(iii)\implies (i) :}$} Since every $SUPM$ is a $UPM$, so $(iii)\implies (i)$.\\
%This completes the proof.
{\bf Proof of Theorem 1.2 :} Let us consider $\mathbb{F}$ and $\mathbb{G}$ as defined by (\ref{el05}). % \bea\label{el4} F= R(f), \;\;\;\; G=R(\mathcal{L}),\eea where $R(z)$ defined by the equation (\ref{el2}).
Let $f$ be a non-constant meromorphic function and $\mathcal{L}$ be a non-constant L-function such that  %$P(f)=P(\mathcal{L})\implies f=\mathcal{L}$ and
$E_f(S,t)=E_{\mathcal{L}}(S,t)$, where $t\in\mathbb{N}\cup\{0\}$. Then $\mathbb{F}$, $\mathbb{G}$ share $(1,t)$. 
Now for  \begin{itemize}
\item [(a)]$t\geq 2$ and $n\geq 2l+5$, or
\item [(b)] $t=1$ and $n\geq 2l+6$, or
\item [(c)] $t=0$ and $n\geq 2l+11$,
\end{itemize}
in view of  {\em Lemma \ref{l1.1}} we have
\bea\label{ep42}\frac{1}{\mathbb{F}-1}=\frac{A}{\mathbb{G}-1}+B,\eea where $A(\neq0), B\in\mathbb{C}$. % and 


From (\ref{ep42})  we easily obtain \bea \label{ep6}T(r,f) = T(r,\mathcal{L}) + S(r,f).\eea
\par Now in view of  $\ol{N}(r,\infty,\mathcal{L})=O(\log r)$, (\ref{ep6}) and from the construction of $\mathbb{F}$ and $\mathbb{G}$ we get 
\beas &&
\ol{N}(r,0; \mathbb{F})+	\ol{N}(r,\infty; \mathbb{F})+	\ol{N}(r,0; \mathbb{G})+	\ol{N}(r,\infty; \mathbb{G})
\\ &\leq&\ol{N}(r,0;f)+\sum_{j=1}^{k}\ol{N}(r,\beta_j;f)+\ol{N}(r,\infty;f)+\ol{N}(r,0;\mathcal{L})+\sum_{j=1}^{k}\ol{N}(r,\beta_j;\mathcal{L})+\ol{N}(r,\infty;\mathcal{L})
\\ &\leq& (2+k)T(r,f)+(1+k)T(r,\mathcal{L})+O(\log r)\\
&=&(2k+3)T(r,f)+O(\log r)\\
&<&\frac{2k+4}{n}T(r,\mathbb{F})\\\nonumber
&\leq&T(r,\mathbb{F}){\;\;}[\because\;\; n\geq 2k+4].
\eeas
So in view of {\em Lemma \ref{l2}}, we have either $\mathbb{F}\mathbb{G} \equiv 1$ or $\mathbb{F} \equiv \mathbb{G}$.
Again in view of {\em Lemma \ref{l04}} we have  $\mathbb{F}\mathbb{G} \not\equiv1$. Thus $\mathbb{F} \equiv \mathbb{G}$; i.e., $R(f)=R(\mathcal{L})$.

Therefore we find that $E_f(S, t)=E_\mathcal{L}(S,t)\implies f=\mathcal{L}$, whenever $R(f)=R(\mathcal{L})\implies f=\mathcal{L}$. That is $(i)\implies(ii).$ 
\vspace{0.1in}\par To show $(ii) \implies (i)$, suppose that $E_f(S, t)=E_\mathcal{L}(S,t)\implies f=\mathcal{L}$.
Let $R(f)=R(\mathcal{L})$, then we have $R(f)-1=R(\mathcal{L})-1$; i.e., $\d\frac{P(f)}{\phi(f)}=\d\frac{P(\mathcal{L})}{\phi(\mathcal{L})}.$ Therefore  $f$ and  $\mathcal{L}$ share $(S,\infty)$, which implies $E_f(S, t)=E_\mathcal{L}(S,t)$ and hence  $f=\mathcal{L}$.
\\\\{\bf Proof of Corollory \ref{c1} and Corollary \ref{c2}:}	The proof of {\em Corollory \ref{c1}} and {\em Corollary \ref{c2}} directly follows from the proof of {\em Theorem \ref{t1}} and {\em Theorem \ref{t2}} respectively, considering $\ol N(r,\infty;f)=O(\log r)$. Hence we omit the proof.
% 	\section{Statements and Declarations}
% 	\noindent {\bf Funding:} There is no funding from any source for this article.\\
%	{\bf Conflict of interest:} Both the authors of this article declare that they do not have any conflict of interest.\\
%	{\bf Ethics approval:} This article does not contain any studies with human participants or animals performed by any of the authors.
\begin{ack}
\par Sanjay Mallick is thankful to``Science and Engineering Research Board, Department of Science and Technology, Government of India"  for financial support to pursue this research work under the Project File No. EEQ/2021/000316.
\par	 Ripan Saha is thankful to the Department of Atomic Energy (DAE), India, for financial support to pursue this work $[$No. $0203/13(41)/2021$-$R\&D$-$II/13168]$. 	
\end{ack}
\begin{thebibliography}{50}
\bibitem{aljahari} T. C. Alzahary, {\em Meromorphic functions with weighted sharing of one set}, Kyungpook Math. J., {\bf 47}(2007), 57-68.
\bibitem{ab1} A. Banerjee,{\em Uniqueness of meromorphic functions sharing two sets with finite weight II}, Tamkang J. Math., {\bf 41}(4)(2010), 379-392.
\bibitem{bl} A. Banerjee, {\em Fujimoto’s theorem - a further study}, J. Contemp. Math. Anal., {\bf 51}(2016), 199-207.
\bibitem{bk} A. Banerjee and A. Kundu,{\em Weighted value sharing and uniqueness problems concerning L-functions and certain meromorphic functions}, Lith. Math. J., {\bf61}(2021), 161–179.
\bibitem{bk1} A. Banerjee and A. Kundu, {\em Uniqueness of L-function with special class of meromorphic function under restricted sharing of sets}, arXiv:2103.00731v1 [math.CV] 1 Mar 2021.
%\bibitem{banerjee2015new} A. Banerjee, A new class  of strong uniqueness polynomials satisfying Fujimoto's condition, Annales Academiae Scientiarum Fennicae MathematicaVolumen 40, 2015, 465-474.
\bibitem{3a} A. Banerjee and S. Mallick, {\em On the characterisations of a new class of
strong uniqueness polynomials generating
unique range sets}, Comput. Methods Funct. Theory, {\bf 17}(2017), 19-45.

\bibitem{frank1998unique} G. Frank and M. Reinders, {\em A unique range set for meromorphic functions with 11 elements}, Complex Var.Theory Appl., {\bf 37}(1)(1998), 185-193.

%\bibitem{fujimoto2000uniqueness} H. Fujimoto, On uniqueness of meromorphic functions sharing finite sets, Amer. J. Math., 122(2000),1175-1203.

\bibitem{fg} F. Gross, {\em Factorization of meromorphic functions and some open problems}, in J.D. Buckholtz and T.J. Suffridge
(Eds.), \emph{Complex Analysis}. Proceedings of the Conference Held at the University of Kentucky, May 18–22, 1976, Lect.
Notes Math., Vol. 599, Springer, Berlin, Heidelberg, 1977, pp. 51-69.

\bibitem{wkh}W. K. Hayman, {\em Meromorphic Functions}, Oxford Mathematical Monographs, Clarendon Press,
Oxford, 1964.

\bibitem{hu2016} P. C. Hu and B. Q. Li, {\em A simple proof and strengthening of a uniqueness theorem for L-functions}, Can. Math. Bull.,
{\bf59}(2016), 119-122.

\bibitem{kac99}
J. Kaczorowski and A. Perelli, {\em On the structure of the Selberg class, I: $0 \leq d \leq 1$}, Acta Math., {\bf182}(1999), 207-241.

%\bibitem{6} I. Lahiri, Value distribution of certain differential polynomials, Int. J. Math. Math. Sci., 28(2)(2001), 83-91.

\bibitem{lahiriNagoya} I. Lahiri, {\em Weighted sharing and uniqueness of meromorphic functions}, Nagoya Math. J., {\bf 161}(2001), 193-206.

\bibitem{Lahiri CVEE} I. Lahiri, {\em Weighted value sharing and uniqueness of meromorphic functions}, Complex Var. Theory Appl., {\bf 46}(2001), 241-253.

%\bibitem{b4} W. C. Lin and H. X. Yi, Uniqueness Theorems for meromorphic functions that share 	three sets, Complex Variables, 48(2003), 315-327.
\bibitem{li2010}
B. Q. Li, {\em A result on value distribution of	L-functions}, Proc. Amer. Math. Soc. {\bf 138}(2010), 2071-2077.

%\bibitem{li-yi} X. M. Li and H. X. Yi, Results on value distribution of L-functions, Math.Nachr., 286(2013), 1326-1336.	
\bibitem{smtbilisi} S. Mallick, {\em On the existence of unique range sets generated by
non-critically injective polynomials and related issues}, Tbilisi Math. J., {\bf 13} (4) (2020), 81-101.
\bibitem{smfil21} S. Mallick, {\em Unique Range Sets - A Further Study}, Filomat, {\bf 34} (5) (2020), 1499-1516.
\bibitem{sh} P. Sahoo and S. Haldar, {\em Results on L functions and certain uniqueness question of gross}, Lithuanian Math. J.,
{\bf 60}(1)(2020), 80-91.
\bibitem{ss} P. Sahoo and A. Sarkar, {\em On the uniqueness of L-function and meromorphic function sharing one or two
sets}, An. Stiint. Univ. Al. I. Cuza Iasi. Mat., {\bf66}(1)(2020), 81-92.
\bibitem{selberg1992} A. Selberg, {\em Old and new conjectures and results about a class of Dirichlet series, in}: Proccedings of the Amalfi
Conference on Analytic Number Theory (Maiori, 1989), Univ. Salerno, Salerno, 1992, pp. 367-385.

\bibitem{steuding2007value} J. Steuding, {\em Value distribution of L-functions}, Lect. Notes Math., Vol. 1877, Springer-Verlag, Berlin, 2007.


\bibitem{hxy0}  H. X. Yi,{\em Meromorphic functions that share one or two values},  Complex Var. Theory Appl., {\bf28}(1995), 1-11.
%\bibitem{yi1996unicity}
%H. X. Yi, Unicity theorems for meromorphic or entire functions III, Bull. Austral. Math. Soc., 53(1996), 71-82.

\bibitem{15} H. X. Yi, {\em Meromorphic functions that share one or two values II}, Kodai Math. J., {\bf 22}(1999), 264-272.
%\bibitem{yilu1} H.X. Yi and W.R. L\"{u}, Meromorphic functions that share two sets II, Acta Math. Sci. Ser. B (Engl. Ed.), 24(1)(2004), 83-90.
\bibitem{yuan2018}
X. M. Li, Q. Q. Yuan and H. X. Yi,{\em Value distribution of L-functions and uniqueness questions of F. Gross}, Lithuanian
Math. J., {\bf 58}(2)(2018), 249-262.
%\bibitem{hxy0}  H. X. Yi, Meromorphic functions that share one or two values,  Complex Var. Theory Appl., 28(1995), 1-11.
%	\bibitem{li1996unique} P. Li and C.C. Yang, \emph{On the unique range set of meromorphic functions}, Proc. Am. Math. Soc., 124:177-185, 1996.

%	\bibitem{li2004meromorphic} X.M. Li and H.X. Yi, \emph{Meromorphic functions sharing three values}, J. Math. Soc. Japan, 56:147-167, 2004.
%	\bibitem{an2004strong} T. T. H. An, J. T. Wang and P. Wong, \emph{Strong uniqueness polynomials: the complex case}, Complex Variables Theory Appl., 49(1)(2004), 25-54.

%	\bibitem{fujimoto2000uniqueness} H. Fujimoto, \emph{On uniqueness of meromorphic functions sharing finite sets}, Amer. J. Math., 122(2000),1175-1203.
%\bibitem{yilu1} H.X. Yi and W.R. L\"{u}, Meromorphic functions that share two sets II, Acta Math. Sci. Ser. B (Engl. Ed.),
%24(1)(2004), 83-90.
%\bibitem{ab1} A.Banerjee, \emph{Uniqueness of meromorphic functions sharing two sets with finite weight II}, Tamkang J. Math., 41(4)(2010), 379-392.

%\bibitem{bl} A. Banerjee, Fujimoto’s theorem - a further study, J. Contemp. Mathemat. Anal., 51(2016), 199-207.


\end{thebibliography}

\end{document}


@author:
@affiliation:
@title:
@language: English
@pages:
@classification1:
@classification2:
@keywords:
@abstract:
%@filename:
%@DOI:
%@editor:
%@reviewer:
@EOI


